\documentclass[11pt]{book}

% ── Packages ─────────────────────────────────────────────────────────────
\usepackage[margin=1.2in]{geometry}
\usepackage{amsmath,amssymb,amsthm}
\usepackage{mathtools}
\usepackage{setspace}
\usepackage{parskip}
\usepackage{titlesec}
\usepackage{enumitem}
\usepackage[colorlinks=true,linkcolor=black,citecolor=black]{hyperref}
\usepackage{tikz}
\usepackage[splitindex]{imakeidx}
\makeindex[name=topics,title=Topic Index,intoc]
\makeindex[name=vocab,title=Vocabulary Index,intoc]
\usetikzlibrary{arrows.meta,positioning,calc,decorations.pathmorphing,patterns,backgrounds}

% ── Unified chapter figure style ─────────────────────────────────────
\tikzset{
  >=Latex,
  chapterbox/.style={draw, thick, rounded corners},
  flow/.style={->, thick},
  thinflow/.style={->, thin},
  point/.style={circle, fill, inner sep=1.2pt},
  every node/.style={font=\small},
  cfig/.style={thick},
  cfigdash/.style={dashed, thick},
  cfigfill/.style={fill opacity=0.75},
  cfignode/.style={circle, fill, inner sep=1.5pt}
}

% ── Typography ───────────────────────────────────────────────────────────
\setstretch{1.15}
\setlength{\parindent}{0pt}
\setlength{\parskip}{6pt}

\titleformat{\chapter}[display]
  {\normalfont\Large\bfseries}{\chaptername\ \thechapter}{10pt}{\huge}
\titleformat{\section}{\normalfont\large\bfseries}{\thesection}{1em}{}
\titleformat{\subsection}{\normalfont\normalsize\bfseries}{\thesubsection}{1em}{}

% ── Theorem environments ─────────────────────────────────────────────────
\theoremstyle{plain}
\newtheorem{theorem}{Theorem}[chapter]
\newtheorem{proposition}[theorem]{Proposition}
\newtheorem{corollary}[theorem]{Corollary}

\theoremstyle{definition}
\newtheorem{definition}[theorem]{Definition}
\newtheorem{example}[theorem]{Example}
\newtheorem{principle}[theorem]{Principle}

\theoremstyle{remark}
\newtheorem{remark}[theorem]{Remark}
\newtheorem{openproblem}[theorem]{Open Problem}
\newtheorem{exercise}[theorem]{Exercise}

% ── Notation ─────────────────────────────────────────────────────────────
\newcommand{\C}{\mathcal{C}}
\newcommand{\A}{\mathcal{A}}
\newcommand{\G}{\Gamma}
\newcommand{\Ss}{S}

% ── Extension environments ────────────────────────────────────────────────
% Historical preamble: opens each chapter with a paradigm's history
\newenvironment{historicalpreamble}
  {\begin{quote}\small}
  {\end{quote}}

% Historical figures: key contributors and sources
\newenvironment{historicalfigures}
  {\begin{quote}\small\textbf{Key figures.}\enspace\ignorespaces}
  {\end{quote}}

% Timeline: paradigm shift chronology
\newenvironment{paradigmtimeline}
  {\begin{quote}\small\textbf{Timeline.}\enspace\ignorespaces}
  {\end{quote}}

% Formal correspondence: precise mathematical restatement of the chapter idea
\newenvironment{formalcorrespondence}
  {\begin{quote}\small\itshape}
  {\end{quote}}

% Computational aside: computational paradigm parallel
\newenvironment{computationalaside}
  {\begin{quote}\small\itshape}
  {\end{quote}}

% Spherepop comparison: BNF fragment + categorical diagram
\newenvironment{spherepopcomparison}
  {\begin{quote}\small}
  {\end{quote}}

% UML / categorical diagram wrapper
\newenvironment{umldiagram}
  {\begin{center}\begin{tikzpicture}[
      box/.style={draw, rounded corners, align=center,
                  minimum width=3.2cm, minimum height=1cm, font=\small},
      arrow/.style={-Latex, thick}
    ]}
  {\end{tikzpicture}\end{center}}

% Teacher note: pedagogical guidance for instructors
\newenvironment{teachernote}
  {\begin{quote}\small\bfseries\textit{Teacher Note.}\normalfont\small\enspace\ignorespaces}
  {\end{quote}}

% Lesson plan: multi-year curriculum block
\newenvironment{lessonplan}
  {\begin{quote}\small}
  {\end{quote}}

% University bridge: connections to university-level topics
\newenvironment{universitybridge}
  {\begin{quote}\small\itshape\textbf{University connections:}\normalfont\small\itshape\enspace\ignorespaces}
  {\end{quote}}

% Reading list: per-chapter further reading
\newenvironment{readinglist}
  {\begin{quote}\small}
  {\end{quote}}

% ── Title ────────────────────────────────────────────────────────────────
\title{\textbf{Structural Thinking}\\[6pt]
\large A History of Computation and the Geometry of Ideas}
\author{Flyxion}
\date{}

% ════════════════════════════════════════════════════════════════════════
\begin{document}
% ════════════════════════════════════════════════════════════════════════

\frontmatter
\maketitle
\tableofcontents

% ════════════════════════════════════════════════════════════════════════
\chapter*{Teacher's Guide}
\addcontentsline{toc}{chapter}{Teacher's Guide}
% ════════════════════════════════════════════════════════════════════════

\section*{Pedagogical Philosophy}

\begin{teachernote}
This textbook is structured around a single underlying principle:
students learn best when they encounter the same structure in multiple
forms before it is formalised.

Each chapter presents a historical problem, an intuitive resolution, a
formal structure, and a computational interpretation. Students are not
expected to master all layers immediately. Understanding accumulates
through repeated exposure across chapters. The goal is not early
precision, but eventual inevitability: the sensation that the answer
could not have been otherwise.

The teacher's role is to hold the structure stable while students find
their footing in each new representation. Resist the temptation to
explain the structure too early. Let students encounter the same idea
in history, in mathematics, and in computation before naming what it is.
\end{teachernote}

\begin{remark}
Each chapter in this book presents the same underlying structure from
five perspectives: (1) historical development of a computational
paradigm, (2) conceptual explanation, (3) formal correspondence,
(4) computational interpretation, and (5) worked example with exercises.
These are not separate topics. They are different descriptions of the
same underlying system. The teacher notes at the end of each chapter
explain how to navigate these layers and what university-level topics
each chapter is preparing students to encounter.
\end{remark}

\section*{Two-Year Lesson Plan}

\begin{lessonplan}
\textbf{Year~1 (Chapters~1--5).}
Focus: constraints, emergence, abstraction, entropy, and temporal
compression. Emphasis on conceptual understanding and worked examples.
Students should finish Year~1 able to describe a solution as the result
of constraint accumulation rather than inspiration.

\textbf{Year~2 (Chapters~6--12).}
Focus: failure modes, consistency, cross-domain mapping, intelligence,
and synthesis. Emphasis on formal reasoning and structured thinking.
Students should finish Year~2 able to recognise the same structure in
different domains and to diagnose failure when global consistency breaks
down.
\end{lessonplan}

\section*{Three-Year Lesson Plan (Recommended)}

\begin{lessonplan}
\textbf{Year~1 (Chapters~1--3).}
Focus: constraints, structural emergence, and abstraction. Goal:
students learn to think in terms of elimination and invariants. By the
end of the year, students should be able to describe any narrowing
process---in mathematics, language, or daily reasoning---as constraint
accumulation.

\textbf{Year~2 (Chapters~4--7).}
Focus: entropy, time and relegation, local versus global consistency,
and obstruction. Goal: students understand limits, structural breakdown,
and the distinction between local correctness and global coherence. They
should encounter, at least once, a system that looks correct everywhere
and fails globally.

\textbf{Year~3 (Chapters~8--12).}
Focus: functorial correspondence, intelligence as constraint maintenance,
the geometry of ideas, deliberate constraint accumulation, and synthesis.
Goal: students integrate the concepts across domains and begin to reason
formally. By the end they should be able to identify when two apparently
different problems share the same underlying structure.
\end{lessonplan}

\section*{Assessment Philosophy}

\begin{teachernote}
Assessment should not prioritise memorisation of definitions. Instead,
evaluate: the ability to recognise structure across different contexts;
the ability to detect inconsistency; and the ability to explain why a
solution must be what it is, rather than merely that it is. Strong
students will begin to describe answers as \emph{forced} rather than
\emph{chosen}. That shift in language is the clearest sign that the
framework has been internalised.
\end{teachernote}

% ════════════════════════════════════════════════════════════════════════
\chapter*{How to Use This Book}
\addcontentsline{toc}{chapter}{How to Use This Book}
% ════════════════════════════════════════════════════════════════════════

Each chapter unfolds in the same sequence. A historical preamble
introduces a computational paradigm and the problem it was invented to
solve. The conceptual core develops the chapter's main idea in plain
terms. A formal correspondence restates that idea in mathematical
language. A computational perspective shows how the same structure
appears in programming. A worked example walks through the idea
concretely. Exercises ask the student to apply it.

These layers are not independent. The history motivates the concept. The
formalism clarifies the concept. The computation instantiates the
concept. The worked example proves that the concept works in practice.
A student who reads only one layer has read one coordinate of the same
point.

Teacher notes appear at the end of each chapter. They are addressed to
instructors and describe common misconceptions, teaching strategies, and
connections to university-level mathematics and computer science.

A student reading this book for the first time should read the
historical preamble and conceptual core of each chapter before
attempting the formal correspondence. The formal correspondence will
then feel like a restatement of something already understood, not a new
idea.

% ════════════════════════════════════════════════════════════════════════
\chapter*{A Note on Computation}
\addcontentsline{toc}{chapter}{A Note on Computation}
% ════════════════════════════════════════════════════════════════════════

Computation is often described as calculation: the production of a
numerical result from numerical inputs. This description is accurate but
incomplete.

At a deeper level, computation is transformation. A computing system
takes a representation of some structure and produces a new
representation according to rules. The rules may be arithmetic, but
they may equally be symbolic, logical, or structural. What matters is
not the content of the representations but the rules that govern their
transformation and the stable forms those transformations eventually
reach.

This view of computation---as rule-governed transformation converging on
stable forms---is what connects the history of computing machines to the
theory of constraints, abstraction, and invariants that this book
develops. Computation and structured thinking are not analogous. They
are instances of the same process.

% ════════════════════════════════════════════════════════════════════════
\chapter*{A Note on Structure}
\addcontentsline{toc}{chapter}{A Note on Structure}
% ════════════════════════════════════════════════════════════════════════

Structure is not the same as content. Two objects can share the same
structure while having entirely different content, just as two sentences
can have the same grammatical form while meaning entirely different
things.

In this book, structure means constraint. A structure is defined by what
it rules out: the configurations, transformations, or combinations that
are not admissible. The more constraints a system satisfies, the more
determinate its structure. A fully determined structure is one where
every admissible configuration is forced by the constraints themselves.

This definition has a counterintuitive implication: adding constraints
does not restrict understanding. It produces it. The accumulation of
constraints is what makes a system legible, a solution inevitable, and
an idea recognisable.

% ════════════════════════════════════════════════════════════════════════
\chapter*{Notation and Prerequisites}
\addcontentsline{toc}{chapter}{Notation and Prerequisites}
% ════════════════════════════════════════════════════════════════════════

This book assumes familiarity with the following informal concepts.

\textbf{Sets.} A set is a collection of objects. We write $x \in A$ to
mean that $x$ is an element of the set $A$, and $A \subseteq B$ to mean
that every element of $A$ is also an element of $B$. The intersection
$A \cap B$ is the set of elements belonging to both $A$ and $B$.

\textbf{Functions.} A function $f : X \to Y$ assigns to each element of
$X$ exactly one element of $Y$. We call $X$ the domain and $Y$ the
codomain.

\textbf{Logarithms.} The logarithm $\log_b n$ is the exponent to which
$b$ must be raised to produce $n$. We write $\log$ for $\log_2$
throughout, following the convention of information theory\index[topics]{information theory}.

\textbf{Graphs (informal).} A graph is a collection of nodes connected
by edges. Graphs are used informally in this book to depict
relationships between systems and components.

No calculus, linear algebra, or formal proof techniques are required.
Where proofs appear, they are sketched rather than presented in full
technical detail. The mathematical appendix at the back of the book
provides additional background.

% ════════════════════════════════════════════════════════════════════════
\chapter*{Preface}
\addcontentsline{toc}{chapter}{Preface}
% ════════════════════════════════════════════════════════════════════════

This book takes a single idea seriously: that ideas themselves obey
structural laws.

It argues that great insights are not accidents of inspiration but
convergence phenomena---the inevitable endpoint of accumulating enough
constraints from enough different domains. It formalises what that means
using the beginnings of several mathematical traditions: set theory,
logic, topology, information theory, and category theory\index[topics]{category theory}.

No advanced mathematical background is assumed. What is assumed is
willingness to think carefully and tolerance for precise language. The
formalism here is a tool for clarity, not an obstacle to it.

The payoff is practical. Once you understand how constraint accumulation
works---formally---you can deliberately accelerate it.

\begin{remark}
Throughout this text, each chapter opens with a historical account of a
computing paradigm and closes with three parallel layers: a formal
correspondence that restates the chapter's main idea in mathematical
language, a computational aside that shows the same idea as it appears
in programming, and a brief comparison with Spherepop Calculus\index[topics]{Spherepop Calculus}, a
formal system in which structure is built entirely from nested regions,
merges, and collapses. These layers are not supplements to the main
argument. They are the same argument, stated in different coordinate
systems. The goal is for each restatement to feel like a clarification
of something already understood, not a new idea.
\end{remark}

\mainmatter

% ════════════════════════════════════════════════════════════════════════
% Global Structural Note
% ════════════════════════════════════════════════════════════════════════

\begin{remark}
Each chapter in this book presents the same underlying structure from
multiple perspectives: (1)~historical development of a computational
paradigm, (2)~conceptual explanation, (3)~formal correspondence,
(4)~computational interpretation, and (5)~worked example and exercises.
These are not separate topics but different descriptions of the same
underlying system. A student who feels that each new layer is restating
something already known is reading correctly.
\end{remark}

% ════════════════════════════════════════════════════════════════════════
\part{Constraints and Their Accumulation}
% ════════════════════════════════════════════════════════════════════════

\chapter{Constraints as the Engine of Thought}

\begin{historicalpreamble}
Early computing machines were not designed to ``find answers.'' They were
designed to follow instructions. From the mechanical calculators of the
nineteenth century to the earliest stored-program computers, computation
was understood as a sequence of explicit steps carried out in order.
Programs were written as procedures: do this, then this, then this. Each
step restricted what could happen next. Errors did not arise from mystery,
but from incorrect instructions or missing conditions.

Over time it became clear that what mattered was not the individual steps,
but the way they progressively eliminated possibilities. A program did not
construct an answer out of nothing---it reduced a space of possibilities
until only one remained. This view of computation as sequential constraint
application forms the foundation of imperative programming\index[topics]{imperative programming} and, as this
chapter argues, the foundation of thought itself.
\end{historicalpreamble}

\begin{historicalfigures}
Charles Babbage\index[topics]{Babbage, Charles} (1791--1871) designed the Analytical Engine, an early
mechanical model of programmable computation. Ada Lovelace\index[topics]{Lovelace, Ada} (1815--1852)
recognised that such a machine could follow symbolic rules, not just
perform numerical calculation. John von Neumann\index[topics]{von Neumann, John} (1903--1957) formalised
the stored-program architecture, establishing sequential execution as the
dominant computational model of the twentieth century.
\end{historicalfigures}

\begin{paradigmtimeline}
1830s: mechanical computation (Babbage). \quad
1940s: stored-program computers. \quad
1950s--60s: assembly languages and early imperative languages. \quad
1970s: structured programming and control-flow discipline.
\end{paradigmtimeline}

% ── Chapter Figure ────────────────────────────────────────────────────
\begin{figure}[h]
\centering
\begin{tikzpicture}[scale=0.92]
  % Outer cube (projected)
  \draw[cfig] (0,0) -- (4,0) -- (4,4) -- (0,4) -- cycle;
  \draw[cfig] (1.2,1.2) -- (5.2,1.2) -- (5.2,5.2) -- (1.2,5.2) -- cycle;
  \draw[cfig] (0,0) -- (1.2,1.2);
  \draw[cfig] (4,0) -- (5.2,1.2);
  \draw[cfig] (4,4) -- (5.2,5.2);
  \draw[cfig] (0,4) -- (1.2,5.2);
  % Constraint planes
  \fill[blue!15,cfigfill] (1.4,0.7) -- (4.6,0.7) -- (4.6,4.7) -- (1.4,4.7) -- cycle;
  \fill[red!15,cfigfill]  (0.8,1.4) -- (4.8,1.4) -- (4.8,2.6) -- (0.8,2.6) -- cycle;
  \fill[green!18,cfigfill](2.2,0.4) -- (3.2,0.4) -- (4.4,1.6) -- (3.4,1.6) -- cycle;
  % Admissible region
  \fill[orange!55,cfigfill] (2.45,1.55) rectangle (3.25,2.45);
  \draw[cfig,orange!70!black] (2.45,1.55) rectangle (3.25,2.45);
  % Labels
  \node at (2.0,-0.55) {configuration space $\Gamma$};
  \node[blue!65!black]   at (5.55,4.4) {$C_1$};
  \node[red!65!black]    at (5.55,2.0) {$C_2$};
  \node[green!55!black]  at (4.5,0.7)  {$C_3$};
  \node[orange!80!black] at (2.85,2.9) {$\mathcal{A}(\mathcal{C})$};
\end{tikzpicture}
\caption{Constraints $C_1, C_2, C_3$ each define a subset of configuration space.
Their intersection is the admissible set $\mathcal{A}(\mathcal{C})$: the only region
that satisfies all conditions simultaneously.}
\end{figure}

\section{Formal Definition}

\begin{definition}[Constraint]\index[vocab]{constraint}
A \textbf{constraint} is a predicate $c : X \to \{\mathrm{true}, \mathrm{false}\}$
defined on a set $X$ of possible configurations. A constraint \emph{eliminates}
from consideration every $x \in X$ for which $c(x) = \mathrm{false}$.
\end{definition}

The set $X$ is the space of all initially possible configurations---
answers, solutions, models. A constraint reduces this space.

\begin{definition}[Constraint Set and Admissible Set]\index[vocab]{admissible set}\index[vocab]{constraint set}
Given a set of constraints $\C = \{c_1, c_2, \ldots, c_n\}$ on $X$, the
\textbf{admissible set} is
\[
  \A(\C) \;=\; \bigl\{ x \in X \;\mid\; c_i(x) = \mathrm{true}
  \text{ for all } i \bigr\}.
\]
\end{definition}

\begin{example}
Let $X = \{1, 2, \ldots, 100\}$. Define constraints:
\begin{align*}
  c_1(x) &= [x \text{ is even}] \\
  c_2(x) &= [x > 50] \\
  c_3(x) &= [x < 70]
\end{align*}
Then $\A(\{c_1, c_2, c_3\}) = \{52, 54, 56, 58, 60, 62, 64, 66, 68\}$.
Nine possibilities remain from one hundred.
\end{example}

As constraints are added, $|\A(\C)|$ typically decreases. When
$|\A(\C)| = 1$, a unique solution is determined. This is the formal
structure of discovery: the constraints forced the answer.

\section{Domains as Constraint Systems}

\begin{definition}[Domain]\index[vocab]{domain}
A \textbf{domain} $D$ is a pair $(X_D, \C_D)$ where $X_D$ is a space of
domain-specific configurations and $\C_D$ is a set of constraints
governing admissibility within $D$.
\end{definition}

Physics, mathematics, music theory, and grammar are all domains in this
sense. Each has its own space of configurations and its own constraints on
what is admissible.

\begin{remark}
A domain expert is someone who has internalized $\C_D$ so thoroughly that
it operates automatically---as a filter rather than a conscious checklist.
This is what aspect relegation formalizes (Chapter~\ref{ch:relegation}).
\end{remark}

\section{Cross-Domain Constraint Accumulation}

The most powerful constraints arise when we require a configuration to be
admissible in multiple domains simultaneously.

\begin{definition}[Cross-Domain Admissibility]\index[vocab]{cross-domain admissibility}
Given domains $D_1 = (X_1, \C_1)$ and $D_2 = (X_2, \C_2)$ with a common
mapping $\phi : X \to X_1 \times X_2$, the \textbf{cross-domain admissible
set} is
\[
  \A(\C_1 \cup \C_2) \;=\; \bigl\{ x \in X \;\mid\;
  \phi(x) \in \A(\C_1) \times \A(\C_2) \bigr\}.
\]
\end{definition}

Each additional domain reduces $|\A|$. When enough domains are imposed,
the admissible set may collapse to a single element.

\begin{principle}[Inevitability Under Density]\index[topics]{inevitability under density}
If the intersection of admissible sets across sufficiently many domains
contains exactly one element, that element is determined by the constraints
alone, not by guessing or inspiration.
\end{principle}

% ── Formal Correspondence ─────────────────────────────────────────────
\section*{Formal Correspondence}
\begin{formalcorrespondence}
Let $\Gamma$ denote the set of all possible configurations of a system and
let $\mathcal{C} = \{C_1, C_2, \dots, C_n\}$ be a set of constraints. Each
$C_i$ defines a subset $\Gamma_i \subseteq \Gamma$. The admissible set is
\[
  \Gamma_{\mathcal{C}} = \bigcap_{i=1}^n \Gamma_i.
\]
An idea corresponds to an element $x \in \Gamma_{\mathcal{C}}$. When
$|\Gamma_{\mathcal{C}}| = 1$, the solution is uniquely determined by the
constraints. This is the same structure stated above, expressed in
set-theoretic terms.
\end{formalcorrespondence}

% ── Computational Aside ───────────────────────────────────────────────
\section*{Computational Perspective}
\begin{computationalaside}
In imperative programming, a system is transformed step by step. Each
instruction reduces what remains possible. A constraint behaves the same
way: it removes configurations from consideration. A sequence of
constraints is therefore not a collection of ideas but a process of
elimination. What appears as a final answer is often just the state that
remains after all eliminations have been applied.
\end{computationalaside}

% ── Worked Example ────────────────────────────────────────────────────
\section*{Worked Example}
Consider the problem of identifying a number between 1 and 100. Impose
three constraints: the number is even; the number is greater than 50; the
number is divisible by 5. Each constraint reduces the admissible set:
\[
  \Gamma_1 = \{2,4,\dots,100\}, \quad
  \Gamma_2 = \{52,54,\dots,100\}, \quad
  \Gamma_3 = \{50,55,60,\dots,100\}.
\]
Their intersection is $\{60,70,80,90,100\}$. Adding a fourth
constraint---the number is less than 80---yields $\{60,70\}$. The solution
is not constructed; it is progressively revealed by eliminating
incompatible possibilities.

\begin{exercise}
Write a sequence of yes/no questions that narrows any number between 1 and
100 to a single value. Observe how each question eliminates possibilities
rather than constructing the answer directly.
\end{exercise}

% ── Spherepop Comparison ──────────────────────────────────────────────
\section*{Spherepop Comparison}
\begin{spherepopcomparison}
Spherepop begins from the idea that structure is built by nested regions,
merges, and collapses. In this chapter, a constraint behaves like a rule
that prunes possible pops until only admissible ones remain.

\textbf{BNF fragment.}
\begin{verbatim}
<system>      ::= <region>
<region>      ::= <atom> | "(" <region-list> ")"
<region-list> ::= <region> | <region> <region-list>
<constraint>  ::= "allow" <region> | "forbid" <region>
\end{verbatim}

A constraint-preserving system retains only regions that are not forbidden.
The admissible set corresponds exactly to the pops that survive elimination.
This expresses the same structure introduced in the definitions above, now
stated as a formal grammar.
\end{spherepopcomparison}

\begin{umldiagram}
  \node[box] (space)  {Configuration\\Space};
  \node[box, right=2.4cm of space]  (filter) {Constraint\\Filter};
  \node[box, right=2.4cm of filter] (admit)  {Admissible\\Regions};
  \draw[arrow] (space) -- (filter);
  \draw[arrow] (filter) -- (admit);
\end{umldiagram}

% ── Teacher Notes ─────────────────────────────────────────────────────
\section*{Teacher Notes}

\begin{teachernote}
Students should begin to see thinking as elimination rather than
construction. The most common difficulty is that students try to guess
answers rather than narrow possibilities systematically. Force step-by-step
elimination exercises and do not permit jumping ahead. Watch for students
who confuse rules with outcomes---who treat a constraint as a description
of the answer rather than as a filter on possibilities.

The worked example in this chapter is deliberately numerical. Use it to
establish the habit of making the admissible set explicit at each step.
Later chapters depend on students having internalised this habit.
\end{teachernote}

\begin{universitybridge}
Constraint satisfaction problems (CSPs) in artificial intelligence.
Linear programming and feasible region analysis. Propositional logic
and satisfiability (SAT). Integer programming.
\end{universitybridge}

\section*{Further Reading}

\begin{readinglist}
George P\'olya, \textit{How to Solve It} --- an introduction to
problem solving as a process of narrowing possibilities; accessible at
any level.

Thomas H. Cormen et al., \textit{Introduction to Algorithms} (selected
sections) --- shows how constraints shape algorithm design.

Rina Dechter, \textit{Constraint Processing} --- a formal treatment of
constraint satisfaction problems at university level.
\end{readinglist}

\chapter{The Three Conditions for Structural Emergence}

\begin{historicalpreamble}
In the early twentieth century, mathematicians sought to understand what it
meant to compute in the most abstract sense. Alonzo Church introduced the
lambda calculus\index[topics]{lambda calculus}, a system in which computation was no longer described as
a sequence of steps but as the transformation of expressions. In this
framework, one does not execute instructions; one reduces an expression
according to formal rules until it reaches a stable form. This stable
form---called a normal form\index[topics]{normal form}---represents the completed result of the
computation.

This shift was profound. It suggested that computation is not inherently
procedural. It can instead be understood as the process of eliminating
instability within a structure until no further change is possible. The
same logic applies to ideas: an idea is not assembled but simplified to
the point where it can no longer be altered without violating a constraint.
Modern functional programming\index[topics]{functional programming} languages inherit this view, treating programs
as expressions to be evaluated rather than instructions to be followed.
\end{historicalpreamble}

\begin{historicalfigures}
Alonzo Church\index[topics]{Church, Alonzo} (1903--1995) developed lambda calculus as a formal system
of computation. Alan Turing\index[topics]{Turing, Alan} (1912--1954) introduced the Turing machine
and proved the two models equivalent. Haskell Curry\index[topics]{Curry, Haskell} (1900--1982) advanced
combinatory logic, which directly influenced the design of functional
programming languages.
\end{historicalfigures}

\begin{paradigmtimeline}
1930s: lambda calculus and formal models of computation. \quad
1950s: theoretical equivalence of computation models proven. \quad
1980s--90s: functional languages (ML, Haskell). \quad
2000s: functional paradigms integrated into mainstream languages.
\end{paradigmtimeline}

% ── Chapter Figure ────────────────────────────────────────────────────
\begin{figure}[h]
\centering
\begin{tikzpicture}[scale=0.98]
  \fill[blue!18,cfigfill]   (0,0)   ellipse (2.1 and 1.3);
  \fill[red!18,cfigfill]    (1.9,0) ellipse (2.1 and 1.3);
  \fill[green!20,cfigfill]  (0.95,1.35) ellipse (2.1 and 1.3);
  \draw[cfig,blue!55!black]  (0,0)   ellipse (2.1 and 1.3);
  \draw[cfig,red!55!black]   (1.9,0) ellipse (2.1 and 1.3);
  \draw[cfig,green!45!black] (0.95,1.35) ellipse (2.1 and 1.3);
  % Overlap region
  \fill[orange!60,cfigfill] (0.95,0.42) circle (0.44);
  \draw[cfig,orange!80!black] (0.95,0.42) circle (0.44);
  % Labels
  \node[blue!65!black]    at (-1.5,-1.6) {constraint density};
  \node[red!65!black]     at ( 3.4,-1.6) {closure};
  \node[green!50!black]   at ( 0.95, 2.9) {projection capacity};
  \node[orange!85!black]  at ( 0.95, 0.42) {emergence};
\end{tikzpicture}
\caption{Structural emergence appears only at the intersection of all three conditions:
constraint density, closure, and projection capacity. Each condition alone is necessary
but not sufficient.}
\end{figure}

\begin{definition}[Constraint Density]\index[vocab]{constraint density}
A system is \textbf{constraint-dense} if $|\A(\C)|$ is small relative to
$|X|$ and the constraints interact non-trivially---that is, $\A(\C_1 \cup
\C_2) \subsetneq \A(\C_1) \cap \A(\C_2)$ for many pairs $C_1, C_2 \subset \C$.
\end{definition}

\begin{definition}[Closure]\index[vocab]{closure}
A system satisfies the \textbf{closure condition} if $\A(\C) \neq \emptyset$.
Closure is the precondition for any solution to exist at all.
\end{definition}

\begin{definition}[Projection Capacity]\index[vocab]{projection capacity}
A system has \textbf{projection capacity} if there exists a map
$\Pi : X \to Y$ into a simpler space $Y$ such that $\Pi$ preserves the
essential structure of $\A(\C)$: the admissible configurations in $Y$
correspond to those in $X$ under the same constraints.
\end{definition}

\begin{theorem}[Necessary and Sufficient Conditions]\index[topics]{structural emergence, conditions for}
Cross-domain structural ideas emerge if and only if the generating system
is constraint-dense, satisfies closure, and possesses projection capacity.
\end{theorem}

\begin{proof}[Sketch]
Constraint density ensures that constraints from different domains interact,
ruling out configurations that would survive each domain individually.
Closure ensures a solution exists. Projection capacity ensures the solution
can be represented in a compressed form that preserves the relevant
structure. Without any one of these, either no solution is determined, no
solution exists, or the solution cannot be expressed.
\end{proof}

\begin{remark}
This theorem reframes the question of intellectual originality. Rather than
asking ``how did this person come up with this?'' we ask: ``was the system
they operated within constraint-dense, closed, and projection-capable?'' If
yes, the idea was in some sense inevitable.
\end{remark}

% ── Formal Correspondence ─────────────────────────────────────────────
\section*{Formal Correspondence}
\begin{formalcorrespondence}
Let $T$ be a transformation that applies constraints to a candidate
configuration. A stable configuration satisfies
\[
  T(x) = x.
\]
Such configurations are fixed point\index[topics]{fixed point}s. In computational systems---for
instance, lambda calculus---these correspond to normal forms: structures
that no longer change under evaluation. Recognition is the process of
reaching such a fixed point. Constraint density, closure, and projection
capacity are exactly the conditions that guarantee a fixed point exists,
is unique, and is representable. This expresses the same structure as
the theorem above, now stated in the language of iteration.
\end{formalcorrespondence}

% ── Computational Aside ───────────────────────────────────────────────
\section*{Computational Perspective}
\begin{computationalaside}
In functional programming, one does not construct results step by step.
Instead, one defines expressions and allows them to reduce. An idea
behaves similarly: it is not assembled but simplified. The final form
is reached when no further transformation changes it. Recognition is the
moment at which reduction has completed.
\end{computationalaside}

% ── Worked Example ────────────────────────────────────────────────────
\section*{Worked Example}
Consider a system defined by three constraints:
\[
  C_1: x + y = 10, \qquad C_2: x - y = 2, \qquad C_3: x,y \in \mathbb{Z}.
\]
Adding $C_1$ and $C_2$ gives $2x = 12$, so $x = 6$. Substituting into
$C_1$ gives $y = 4$. The system stabilises at $(x,y) = (6,4)$.
Structure emerges only when all three constraints are satisfied
simultaneously; removing any one allows multiple solutions.

\begin{exercise}
Simplify the expression $((x + 0) \times 1)$ until no further simplification
is possible. Notice that the result is not constructed but revealed by
removing what does not change the value.
\end{exercise}

% ── Spherepop Comparison ──────────────────────────────────────────────
\section*{Spherepop Comparison}
\begin{spherepopcomparison}
Spherepop reduction reaches closure when no further collapse is required.
This matches the chapter's account of structural emergence.

\textbf{BNF fragment.}
\begin{verbatim}
<expression> ::= <region> | <merge> | <collapse>
<merge>      ::= <region> "+" <region>
<collapse>   ::= "[" <expression> "]" "=>" <region>
\end{verbatim}

A structurally emergent object is one that has passed through enough
merges and collapses to reach a stable region. Constraint density
corresponds to many possible collapse pressures; closure to the
existence of a stable reduced region; projection capacity to whether
the reduced region can serve as a usable outer form. This is the same
structure introduced in the theorem above, now stated as a grammar.
\end{spherepopcomparison}

\begin{umldiagram}
  \node[box] (dense)   {Constraint\\Density};
  \node[box, right=2.4cm of dense]   (closure) {Closure};
  \node[box, right=2.4cm of closure] (stable)  {Stable\\Form};
  \draw[arrow] (dense) -- (closure);
  \draw[arrow] (closure) -- (stable);
\end{umldiagram}

% ── Teacher Notes ─────────────────────────────────────────────────────
\section*{Teacher Notes}

\begin{teachernote}
Students should recognise that solutions appear when constraints interact,
not when effort increases. The most common difficulty is a persistent
belief that insight is random or purely intuitive. Counter this by using
systems where removing one constraint breaks uniqueness: show the student
a system with three constraints that has a unique solution, then remove
one, and ask what changes. Watch for students who solve procedurally
without recognising the structural conditions that make a unique solution
possible.

The three conditions---constraint density, closure, and projection
capacity---should be treated as a checklist students can apply to any
problem. Ask them to identify which condition is violated when a system
fails to produce a solution.
\end{teachernote}

\begin{universitybridge}
Fixed point theory and Banach's fixed point theorem. Lambda calculus and
normal forms. Dynamical systems and convergence. Nonlinear systems and
equilibria.
\end{universitybridge}

\section*{Further Reading}

\begin{readinglist}
Douglas Hofstadter, \textit{G\"odel, Escher, Bach} --- explores how
structure emerges from interacting rules; accessible and motivating.

Alonzo Church, \textit{An Unsolvable Problem of Elementary Number
Theory} (1936) --- introduces lambda calculus; readable with care.

Henk Barendregt, \textit{The Lambda Calculus} --- a deeper formal
treatment of reduction and fixed points at university level.
\end{readinglist}

% ════════════════════════════════════════════════════════════════════════
\part{Compression, Abstraction, and Representation}
% ════════════════════════════════════════════════════════════════════════

\chapter{Abstraction as Constraint-Preserving Reduction}

\begin{historicalpreamble}
As programming languages evolved, it became clear that not all
computations need to be performed immediately. In some systems, values are
computed only when they are required---a strategy known as lazy evaluation\index[topics]{lazy evaluation}.
This idea emerged from both practical and theoretical concerns. In infinite
data structures such as streams, it is impossible to compute everything in
advance; computation must be deferred until a specific portion is needed.
Languages such as Haskell formalised this approach, allowing expressions to
remain unevaluated until their results are demanded.

This revealed something deeper: evaluation is not a single event but a
process that can be partially completed and resumed later. Abstraction
is the result of completing enough of that process that the internal
structure no longer needs to be visited again.
\end{historicalpreamble}

\begin{historicalfigures}
John Backus (1924--2007) developed FORTRAN and later articulated a
functional programming style. Paul Hudak (1952--2015) contributed
centrally to Haskell and the theory of lazy evaluation. Philip Wadler
(1956--) advanced functional programming and formalised evaluation
strategies that underpin modern abstract computation.
\end{historicalfigures}

\begin{paradigmtimeline}
1950s: early compilers (FORTRAN). \quad
1970s: demand-driven evaluation concepts proposed. \quad
1980s--90s: lazy evaluation formalised (Haskell). \quad
2000s: streaming and deferred computation in production systems.
\end{paradigmtimeline}

% ── Chapter Figure ────────────────────────────────────────────────────
\begin{figure}[h]
\centering
\begin{tikzpicture}[scale=0.9]
  % Detailed structure (left)
  \draw[cfig] (0,0) -- (1,1.5) -- (2.2,1.1) -- (2.8,2.2)
              -- (4,1.8) -- (4.5,0.7) -- (3.5,-0.4)
              -- (2.1,-0.7) -- (1.1,-0.2) -- cycle;
  \draw[thin,gray!60] (1,1.5) -- (3.5,-0.4);
  \draw[thin,gray!60] (2.2,1.1) -- (1.1,-0.2);
  \draw[thin,gray!60] (2.8,2.2) -- (2.1,-0.7);
  \node at (2.2,-1.35) {detailed structure $X$};
  % Invariant labels
  \node at (0.65,1.9)  {$a$};
  \node at (4.15,2.15) {$b$};
  \node at (3.8,-0.85) {$c$};
  % Projection arrow
  \draw[->,very thick] (5.2,0.75) -- (7.0,0.75);
  \node at (6.1,1.15) {$\pi$};
  % Abstract structure (right)
  \fill[gray!10] (7.5,0) rectangle (10.5,2.1);
  \draw[cfig]   (7.5,0) rectangle (10.5,2.1);
  \node at (9.0,-0.5) {abstract structure $\tilde{X}$};
  % Invariant labels preserved
  \node at (7.85,2.45)  {$a$};
  \node at (10.15,2.45) {$b$};
  \node at (10.15,-0.3) {$c$};
\end{tikzpicture}
\caption{Abstraction $\pi$ maps detailed internal structure to a simpler representation.
The invariants $a$, $b$, $c$ are preserved; internal distinctions irrelevant to them are discarded.}
\end{figure}

\section{The Formal Definition}

\begin{definition}[Abstraction Operator]\index[vocab]{abstraction operator}\index[vocab]{abstraction}
An \textbf{abstraction operator} is a surjective map $\alpha : X \to Y$
from a high-dimensional configuration space $X$ to a lower-dimensional
space $Y$ such that:
\begin{enumerate}[label=(\roman*)]
  \item $\alpha$ preserves admissibility: if $x \in \A(\C)$ then
        $\alpha(x) \in \A(\alpha(\C))$, where $\alpha(\C)$ denotes the
        image of the constraint set;
  \item $\alpha$ does not identify inadmissible with admissible:
        if $\alpha(x_1) = \alpha(x_2)$ then $x_1 \sim_\C x_2$
        (they are constraint-equivalent).
\end{enumerate}
\end{definition}

Condition (i) says the abstraction does not discard relevant structure.
Condition (ii) says the abstraction only identifies things that are
genuinely equivalent under the constraints.

\section{Information Loss and Structural Preservation}

Every abstraction involves loss. The question is which loss is acceptable.

\begin{principle}[Admissible Loss]\index[topics]{admissible loss}
The information discarded by a well-formed abstraction is exactly the
information that does not affect constraint satisfaction. Its loss
simplifies without distorting.
\end{principle}

\begin{example}
The ideal gas law $PV = nRT$ abstracts over the positions and velocities
of individual molecules. It discards enormous detail. But for the purposes
of predicting bulk behavior under the relevant constraints, the discarded
detail is irrelevant. The abstraction preserves what matters.
\end{example}

\section{Invariants Under Abstraction}

\begin{definition}[Invariant]\index[vocab]{invariant}
A property $P$ of configurations in $X$ is an \textbf{invariant under
abstraction $\alpha$} if $P(x) = P(\alpha(x))$ for all $x \in X$---that
is, the property is preserved by the abstraction.
\end{definition}

The goal of abstraction is to find representations in which invariants
become explicit. A good abstraction makes the structure of the constraint
system legible.

% ── Formal Correspondence ─────────────────────────────────────────────
\section*{Formal Correspondence}
\begin{formalcorrespondence}
Abstraction is modelled as a projection $\pi : X \to \tilde{X}$. Define
an equivalence relation:
\[
  x_1 \sim x_2 \iff \pi(x_1) = \pi(x_2).
\]
Then $\tilde{X} = X / {\sim}$. Abstraction corresponds to quotienting:
internal differences are removed once they no longer affect external
interaction. This aligns with reduction---internal computation is
completed before projection. This expresses the same structure as the
abstraction operator defined above, now stated as a quotient space\index[topics]{quotient space}.
\end{formalcorrespondence}

% ── Computational Aside ───────────────────────────────────────────────
\section*{Computational Perspective}
\begin{computationalaside}
In modular systems, complexity is hidden behind stable interfaces. A
component may contain many internal operations, but externally it behaves
as a single unit. Abstraction follows the same principle: internal
structure is resolved to the point that it no longer affects interaction.
To use an abstraction is to rely on completed computation without
revisiting it.
\end{computationalaside}

% ── Worked Example ────────────────────────────────────────────────────
\section*{Worked Example}
Consider the expression $(3 + 5) \times (2 + 4)$. Expanding fully gives
$8 \times 6 = 48$. Alternatively, set $a = 3 + 5$ and $b = 2 + 4$; then
$a \times b = 48$. The abstraction removes intermediate computation steps
while preserving the final relationship. Different internal structures
collapse to the same result: $(4 + 4)(1 + 5) = 8 \times 6 = 48$.
Abstraction identifies equivalence classes of computations sharing the
same external behaviour.

\begin{exercise}
Describe how you would explain a calculator to someone without describing
its internal circuitry. What information is necessary, and what can be
hidden?
\end{exercise}

% ── Spherepop Comparison ──────────────────────────────────────────────
\section*{Spherepop Comparison}
\begin{spherepopcomparison}
In Spherepop, abstraction means that an inner region has been collapsed
enough that the outer system can treat it as a single unit.

\textbf{BNF fragment.}
\begin{verbatim}
<region>   ::= <atom> | "(" <region-list> ")"
<abstract> ::= "abs(" <region> ")"
<equiv>    ::= <region> "~" <region>
\end{verbatim}

If two regions reduce to the same outer behaviour, they may be treated as
equivalent under abstraction. Replacing a complex nested pop by a single
abstract region while preserving allowed outer interactions is the same
operation as the quotient map $\pi$ above, now stated as a grammar rule.
\end{spherepopcomparison}

\begin{umldiagram}
  \node[box] (inner)   {Nested\\Region};
  \node[box, right=2.4cm of inner]   (collapse) {Abstraction\\(Quotient)};
  \node[box, right=2.4cm of collapse] (outer)   {Abstract\\Unit};
  \draw[arrow] (inner) -- (collapse);
  \draw[arrow] (collapse) -- (outer);
\end{umldiagram}

% ── Teacher Notes ─────────────────────────────────────────────────────
\section*{Teacher Notes}

\begin{teachernote}
Students should understand that abstraction hides completed work, not
incomplete work. The most common misconception is treating abstraction
as simplification---a removal of complexity for convenience---rather than
as the preservation of what matters once internal distinctions have been
resolved. Show different internal forms that produce identical outputs,
and ask students to identify what was preserved and what was discarded.

Watch for students who lose track of what must remain invariant. Ask
explicitly: if I change the internal structure, what would have to
change in the output for this to count as a different abstraction?
\end{teachernote}

\begin{universitybridge}
Equivalence relations and quotient spaces. Type systems and abstract
data types. Category theory: objects defined by morphisms rather than
internal construction. Homotopy type theory.
\end{universitybridge}

\section*{Further Reading}

\begin{readinglist}
Harold Abelson and Gerald Jay Sussman, \textit{Structure and
Interpretation of Computer Programs} --- shows abstraction as the
central tool of computation; freely available online.

Barbara Liskov, \textit{Data Abstraction and Hierarchy} (1987) ---
introduces abstraction in system design; concise and precise.

F. William Lawvere and Stephen Schanuel, \textit{Conceptual
Mathematics} --- accessible introduction to category theory.
\end{readinglist}

\chapter{Multiplicity, Entropy, and Compression}

\begin{historicalpreamble}
As software systems grew in size and complexity, it became impossible for
any one person to understand every detail of a program. This led to the
development of modular programming, in which systems are divided into
components with well-defined interfaces. A module hides its internal
workings while exposing only what is needed for interaction, allowing
programmers to build large systems by combining smaller parts without
needing to re-examine their internal structure.

Concepts such as functions, libraries, and application programming
interfaces formalised this approach. A function could be used reliably
once its behaviour was understood, regardless of how it was implemented.
This marked a shift from computation as execution to computation as the
composition of stabilised units---and it is exactly the shift that
multiplicity and entropy make precise.
\end{historicalpreamble}

\begin{historicalfigures}
Edsger Dijkstra\index[topics]{Dijkstra, Edsger} (1930--2002) advocated structured programming and
conceptual clarity. Barbara Liskov\index[topics]{Liskov, Barbara} (1939--) developed foundational
principles of abstraction and data encapsulation. Alan Kay\index[topics]{Kay, Alan} (1940--)
pioneered object-oriented thinking and the idea that meaning lies in
interfaces rather than implementations.
\end{historicalfigures}

\begin{paradigmtimeline}
1960s: structured programming. \quad
1970s: modular systems and data abstraction. \quad
1980s: object-oriented programming\index[topics]{object-oriented programming}. \quad
2000s: APIs and large-scale software ecosystems.
\end{paradigmtimeline}

% ── Chapter Figure ────────────────────────────────────────────────────
\begin{figure}[h]
\centering
\begin{tikzpicture}[scale=0.97]
  % Microstates (left cloud)
  \foreach \x/\y in {0/0, 0.7/1.1, 1.2/0.2, 1.7/1.4, 2.1/0.65,
                     2.4/1.8, 2.8/0.1, 3.0/1.15, 3.3/0.8, 3.6/1.9} {
    \node[cfignode] at (\x,\y) {};
  }
  \node at (1.8,-0.55) {microstates};
  % Flows
  \foreach \x/\y/\tx/\ty in {
    0/0/5.7/0.5,  0.7/1.1/5.7/1.5, 1.2/0.2/5.7/0.5,
    1.7/1.4/5.7/1.5, 2.1/0.65/5.7/0.5, 2.4/1.8/5.7/2.2,
    2.8/0.1/5.7/0.5, 3.0/1.15/5.7/1.5, 3.3/0.8/5.7/1.5,
    3.6/1.9/5.7/2.2} {
    \draw[thinflow,gray!55] (\x,\y) -- (\tx,\ty);
  }
  % Macrostates (right)
  \fill[blue!35]  (5.7,0.15) rectangle (6.4,0.85);
  \fill[red!35]   (5.7,1.15) rectangle (6.4,1.85);
  \fill[green!35] (5.7,1.95) rectangle (6.4,2.55);
  \draw[cfig,blue!60!black]  (5.7,0.15) rectangle (6.4,0.85);
  \draw[cfig,red!60!black]   (5.7,1.15) rectangle (6.4,1.85);
  \draw[cfig,green!50!black] (5.7,1.95) rectangle (6.4,2.55);
  \node at (7.15,0.5)  {$M_1$};
  \node at (7.15,1.5)  {$M_2$};
  \node at (7.15,2.25) {$M_3$};
  \node at (6.05,-0.55) {macrostates};
  % Entropy annotation
  \draw[cfigdash,gray!50] (4.1,0.9) -- (4.1,-0.8);
  \node[gray!60!black] at (4.1,-1.05) {$S = \log\,\mu$};
\end{tikzpicture}
\caption{Compression maps many indistinguishable microstates onto fewer macrostates.
Multiplicity $\mu$ counts how many microstates share a macrostate; entropy $S = \log\,\mu$
measures the uncertainty that compression cannot remove.}
\end{figure}

\section{Multiplicity as a Measure}

\begin{definition}[Multiplicity]\index[vocab]{multiplicity}
Given a constraint set $\C$ on $X$, the \textbf{multiplicity} of a
configuration $y$ in the abstracted space $Y$ is the number of
pre-images in $X$ that map to $y$ under the abstraction:
\[
  \mu(y) \;=\; |\alpha^{-1}(y) \cap \A(\C)|.
\]
\end{definition}

High multiplicity means that many distinct configurations in $X$ are
indistinguishable at the level of $Y$. The abstraction has collapsed them
together.

\section{Entropy as Log-Multiplicity}

\begin{definition}[Structural Entropy]\index[vocab]{entropy!structural}\index[vocab]{entropy}
The \textbf{structural entropy} of a compressed representation is
\[
  \Ss(y) \;=\; \log \mu(y).
\]
\end{definition}

This formalizes the intuition that entropy measures uncertainty or
multiplicity. When $\Ss(y) = 0$, the configuration $y$ corresponds to
exactly one configuration in $X$: the abstraction is lossless at that
point. When $\Ss(y)$ is large, many configurations in $X$ are being
grouped together.

\begin{remark}
Shannon's information entropy \cite{shannon1948} has this same formal
structure. The entropy of a message is the log of the number of messages
that would appear identical under the channel constraints. Our structural
entropy is the same idea applied to constraint systems.
\end{remark}

\section{The Three Roles of Entropy}

In the framework developed here, structural entropy simultaneously plays
three roles that might seem distinct:

\begin{enumerate}
  \item \textbf{Thermodynamic:} In physical systems, it measures the number
        of microscopic configurations compatible with macroscopic observations.
  \item \textbf{Epistemic:} It measures how much of the past cannot be
        recovered from the present---how much historical information is lost
        in compression.
  \item \textbf{Slack:} It measures the residual freedom a system has to
        evolve while remaining within its constraints.
\end{enumerate}

These are not three separate concepts. They are the same concept read at
different scales. This is the kind of unification that abstraction makes
visible.

\section{Constraint-Preserving Compression}

\begin{definition}[Constraint-Preserving Compression]\index[vocab]{compression!constraint-preserving}\index[vocab]{compression}
A compression operator $\Pi : \A(\C) \to \tilde{\A}$ is
\textbf{constraint-preserving} if
\[
  \Pi(\gamma_1) = \Pi(\gamma_2) \;\Rightarrow\;
  \gamma_1 \sim_\C \gamma_2.
\]
That is, $\Pi$ identifies configurations only when they are genuinely
indistinguishable under the constraint set.
\end{definition}

\begin{proposition}[Entropy Under Compression]
Under a constraint-preserving compression $\Pi$,
\[
  \Ss_\Pi(x) \;\leq\; \Ss(x),
\]
with equality if and only if $\Pi$ preserves all admissible distinctions.
\end{proposition}

\begin{proof}
$\Pi$ can only group configurations that are already equivalent under $\C$.
The equivalence classes of $\Pi$ are subsets of the equivalence classes
of $\C$. The multiplicity measure $\mu_\Pi(y)$ is therefore at most
$\mu(y)$ for each $y$, giving $\Ss_\Pi \leq \Ss$.
\end{proof}

% ── Formal Correspondence ─────────────────────────────────────────────
\section*{Formal Correspondence}
\begin{formalcorrespondence}
Abstraction is a projection $\pi : X \to \tilde{X}$ that identifies
configurations equivalent under $\sim$. Multiplicity counts how many
distinct configurations map to the same observable:
\[
  \mu(y) = |\pi^{-1}(y) \cap \A(\C)|.
\]
Entropy is $S(y) = \log \mu(y)$. Compression reduces $\mu$ while
preserving the distinctions that matter. This expresses the same
structure as the proposition above, now reading entropy directly as
log-multiplicity of the projection's fibres.
\end{formalcorrespondence}

% ── Computational Aside ───────────────────────────────────────────────
\section*{Computational Perspective}
\begin{computationalaside}
Spherepop makes multiplicity visible by allowing many internal regions
to reduce to the same observable outer shell. Different inner states may
satisfy $\mathrm{obs}(r_1) = \mathrm{obs}(r_2) = \cdots = \mathrm{obs}(r_n)$.
Multiplicity counts how many distinct Spherepop reductions remain
indistinguishable under the current projection---and this is the direct
analogue of entropy as log-multiplicity.
\end{computationalaside}

% ── Worked Example ────────────────────────────────────────────────────
\section*{Worked Example}
Consider a system with eight possible states $\Gamma = \{s_1,\dots,s_8\}$.
Entropy is $S = \log_2 8 = 3$. Now impose constraints reducing the
admissible set to $\{s_3, s_7\}$. The new entropy is $S = \log_2 2 = 1$.
Compression reduces multiplicity while preserving admissible structure.
A representation that encodes only membership in $\{s_3, s_7\}$ preserves
the constraint-relevant distinction while discarding irrelevant detail.

\begin{exercise}
Describe a complex object using as few words as possible without losing
its essential structure. What information can be removed, and what must
remain?
\end{exercise}

% ── Spherepop Comparison ──────────────────────────────────────────────
\section*{Spherepop Comparison}
\begin{spherepopcomparison}
In Spherepop, multiplicity appears when many internal regions reduce to
the same outer shell.

\textbf{BNF fragment.}
\begin{verbatim}
<state>        ::= <region>
<observable>   ::= "obs(" <state> ")"
<multiplicity> ::= "mult(" <observable> ")"
\end{verbatim}

Different inner states satisfying $\mathrm{obs}(r_i) = \mathrm{obs}(r_j)$
are counted by $\mathrm{mult}$. This is the grammar-level expression of
log-multiplicity as entropy, the same quantity defined in the proposition
above.
\end{spherepopcomparison}

\begin{umldiagram}
  \node[box] (many)    {Many Internal\\States};
  \node[box, right=2.4cm of many]    (obs)  {Shared\\Observable};
  \node[box, right=2.4cm of obs]     (ent)  {Multiplicity /\\Entropy};
  \draw[arrow] (many) -- (obs);
  \draw[arrow] (obs) -- (ent);
\end{umldiagram}

% ── Teacher Notes ─────────────────────────────────────────────────────
\section*{Teacher Notes}

\begin{teachernote}
Students should come to see entropy as counting possibilities, not
measuring disorder. The word ``entropy'' carries misleading connotations
from everyday use. Begin with simple counting: how many configurations
are consistent with what you know? That count is the multiplicity.
Its logarithm is entropy. Establish this concretely before introducing
any formula.

Watch for confusion between compression and loss of meaning. Students
often assume that a shorter representation is a less faithful one.
Show explicitly that constraint-preserving compression discards only
what was already redundant---the meaning is not reduced, only the
representation.
\end{teachernote}

\begin{universitybridge}
Shannon's information entropy. Statistical mechanics and Boltzmann
entropy. Coding theory and Huffman coding. Minimum description length
and Kolmogorov complexity.
\end{universitybridge}

\section*{Further Reading}

\begin{readinglist}
Claude Shannon, \textit{A Mathematical Theory of Communication} (1948)
--- the foundational text; the first few sections are accessible to a
careful high school student.

David J.~C.~MacKay, \textit{Information Theory, Inference, and
Learning Algorithms} --- accessible introduction to entropy and coding;
freely available online.

Thomas Cover and Joy Thomas, \textit{Elements of Information Theory}
--- standard university-level treatment.
\end{readinglist}

% ════════════════════════════════════════════════════════════════════════
\part{Temporal Structure and Relegation}
% ════════════════════════════════════════════════════════════════════════

\chapter{Aspect Relegation and the Arrow of Time}
\label{ch:relegation}

\begin{historicalpreamble}
With the rise of networks and large-scale computing, programs were no
longer confined to a single machine. Distributed systems introduced new
challenges: multiple components operating independently, communicating
across unreliable connections. It became clear that correctness could no
longer be verified locally. Each component might behave correctly on its
own, yet the system as a whole could fail due to mismatched assumptions,
delayed messages, or inconsistent states.

But distributed systems\index[topics]{distributed systems} also revealed something about time. Events on
different machines could not always be placed in a single global order.
Leslie Lamport introduced logical clocks to impose a partial order on
events without requiring a universal clock. This formalised the idea that
time in a computing system is not a background against which computation
happens, but a structure produced by computation itself. Aspect relegation
is the cognitive analogue: the compression of prior computation into a form
that can be retrieved without being reconstructed.
\end{historicalpreamble}

\begin{historicalfigures}
Leslie Lamport\index[topics]{Lamport, Leslie} (1941--) developed the theory of logical clocks and
distributed consistency. His 1978 paper on time, clocks, and the ordering
of events in distributed systems remains foundational. Barbara Liskov
(1939--) contributed consistency models that formalised the trade-offs
between local and global coherence. Eric Brewer\index[topics]{Brewer, Eric} (1955--) formulated the
CAP theorem\index[topics]{CAP theorem}, showing that certain global properties cannot be simultaneously
guaranteed.
\end{historicalfigures}

\begin{paradigmtimeline}
1970s: early networked systems and distributed algorithms. \quad
1978: Lamport's logical clocks paper. \quad
2000s: CAP theorem and large-scale distributed computing. \quad
2010s: cloud infrastructure and global consistency systems.
\end{paradigmtimeline}

% ── Chapter Figure ────────────────────────────────────────────────────
\begin{figure}[h]
\centering
\begin{tikzpicture}[scale=0.97]
  % Winding trajectory
  \draw[very thick,blue!60]
    plot[smooth] coordinates
    {(0,0) (0.8,0.65) (1.3,0.1) (2.0,1.1) (2.8,0.5) (3.6,1.5) (4.4,0.8)};
  \node[cfignode,blue!60] at (0,0) {};
  \node[cfignode,blue!60] at (4.4,0.8) {};
  \node at (2.2,-0.55) {extended computation over time};
  % Compression arrow
  \draw[->,very thick] (5.2,0.75) -- (7.1,0.75);
  \node at (6.15,1.15) {$\tau_t$};
  % Latent block
  \draw[cfig,rounded corners] (7.5,0.2) rectangle (10.1,1.4);
  \fill[gray!12] (7.5,0.2) rectangle (10.1,1.4);
  \node at (8.8,0.8) {latent structure $\Pi(\gamma)$};
  % Time axis
  \draw[->] (-0.3,-1.05) -- (10.3,-1.05) node[right]{$t$};
  % irreversibility annotation
  \draw[cfigdash,red!55] (4.4,0.8) -- (4.4,-0.85);
  \node[red!60!black] at (3.6,-0.7) {information lost};
\end{tikzpicture}
\caption{Temporal compression $\tau_t$ collapses an extended history into a latent structure.
The compression is lossy: the prior trajectory cannot be recovered from the result.
This is the formal content of irreversibility.}
\end{figure}

\section{Temporal Compression}

\begin{definition}[Temporal Compression]\index[vocab]{temporal compression}
The map $\tau_t : X^{(-\infty, t)} \to X(t)$ that sends a history of
configurations to the present state is a \textbf{temporal compression
operator} if it is lossy---that is, if distinct histories can map to
the same present state.
\end{definition}

When $\tau_t$ is lossy, the present state does not uniquely determine
the past. This is the formal definition of irreversibility.

\begin{proposition}[Irreversibility Under Lossy Compression]
If $\tau_t$ is lossy, then $\tau_t^{-1}$ does not exist as a function.
The process is irreversible.
\end{proposition}

\begin{proof}
Surjectivity of $\tau_t$ with non-trivial fibers means multiple distinct
pre-images exist. No right inverse is definable.
\end{proof}

\section{Entropy and History}

The entropy of the present state $X(t)$ measures how many histories are
compatible with it:
\[
  \Ss(X(t)) \;=\; \log |\tau_t^{-1}(X(t))|.
\]

As time advances, more and more distinct histories become consistent with
the present configuration. Entropy increases not because the world becomes
``more disordered'' in some vague sense, but because the present state
becomes consistent with a growing number of possible pasts.

This is the arrow of time\index[topics]{arrow of time}, stated precisely.

\section{Aspect Relegation as Directed Temporal Compression}

\begin{definition}[Aspect Relegation]\index[vocab]{aspect relegation}\index[vocab]{relegation}
\textbf{Aspect relegation} is the process by which a cognitive system
applies temporal compression $\tau_t$ to a constraint structure,
storing the result as latent structure that can be executed without
active reconstruction.
\end{definition}

What we call ``intuition'' is the execution of relegated structure. The
compression happened earlier (through practice and learning); execution
now retrieves the compressed result rather than recomputing it.

\begin{proposition}
Aspect relegation reduces the active computational cost of constraint
satisfaction at execution time, at the cost of requiring prior compression
work.
\end{proposition}

\begin{remark}
This proposition explains why expertise feels effortless and why it takes
time to develop. The effort is not reduced; it is relocated in time.
\end{remark}

\section{Failure Modes}

Relegated structures are built for specific constraint environments. When
the environment changes---when new constraints are introduced that violate
assumptions of the stored compression---relegation fails.

\begin{definition}[Relegation Failure]\index[vocab]{relegation failure}
A relegated structure $\Pi(\gamma)$ undergoes \textbf{relegation failure}
when the environment presents a configuration outside $\A(\C)$ as defined
when $\Pi$ was formed, causing the stored compression to produce an
inadmissible result.
\end{definition}

Relegation failure is not a bug. It is the signal that reactivation of
deliberate constraint-processing is necessary. The error is not in the
relegated structure; it is in applying it past its domain of validity.

% ── Formal Correspondence ─────────────────────────────────────────────
\section*{Formal Correspondence}
\begin{formalcorrespondence}
Let $x(t)$ be a time-dependent state. Temporal compression is the map
$\tau_t : X^{(-\infty,t)} \to X(t)$ that sends a full history to the
present state. When $\tau_t$ is lossy, the present does not uniquely
determine the past. Entropy of the present measures how many histories
are compatible with it: $S(X(t)) = \log|\tau_t^{-1}(X(t))|$. As time
advances, $S$ increases because more distinct histories become consistent
with the current state. This is the arrow of time, stated precisely, and
it is the same structure as the temporal compression operator defined
above.
\end{formalcorrespondence}

% ── Computational Aside ───────────────────────────────────────────────
\section*{Computational Perspective}
\begin{computationalaside}
Compiled systems perform their work in advance, producing a form that can
be executed rapidly. Intuition is analogous: the underlying transformations
have already been carried out, leaving only a direct mapping from input to
result. Speed is not the absence of computation but its prior completion.
What appears as intuition is the execution of relegated structure.
\end{computationalaside}

% ── Worked Example ────────────────────────────────────────────────────
\section*{Worked Example}
Consider solving $(2 + 3) \times 4$. Step by step: $2 + 3 = 5$, then
$5 \times 4 = 20$. Once the intermediate step is completed, the internal
structure $(2+3)$ is no longer needed. We retain only $5 \to 20$. The
earlier computation is relegated---it cannot be reconstructed from the
final result alone. Time introduces irreversibility by compressing prior
structure into a simplified form.

\begin{exercise}
Compare solving a problem step by step with instantly recognising the
answer. What must have happened beforehand for immediate recognition to
be possible?
\end{exercise}

% ── Spherepop Comparison ──────────────────────────────────────────────
\section*{Spherepop Comparison}
\begin{spherepopcomparison}
Aspect relegation in Spherepop appears when previously reduced inner
regions are carried forward as latent units rather than recomputed.

\textbf{BNF fragment.}
\begin{verbatim}
<history>    ::= <state> | <state> "->" <history>
<latent>     ::= "store(" <region> ")"
<reactivate> ::= "use(" <latent> ")"
\end{verbatim}

A latent region is a previously collapsed structure retained for future
reuse. Temporal compression corresponds to storing a reduced Spherepop
unit and deploying it later instead of rebuilding the full nested
computation. This is the same operation as $\tau_t$ above, now stated as
a grammar rule on stored regions.
\end{spherepopcomparison}

\begin{umldiagram}
  \node[box] (comp)   {Expanded\\Computation};
  \node[box, right=2.4cm of comp]   (lat)   {Latent\\Unit};
  \node[box, right=2.4cm of lat]    (reuse) {Reactivated\\Use};
  \draw[arrow] (comp) -- (lat);
  \draw[arrow] (lat) -- (reuse);
\end{umldiagram}

% ── Teacher Notes ─────────────────────────────────────────────────────
\section*{Teacher Notes}

\begin{teachernote}
Students should come to see time as compression of prior structure, not
merely as sequence. The most common difficulty is treating temporal order
as just a list of steps rather than as a process that destroys
information as it advances. Begin with concrete reversible and
irreversible processes: scrambling an egg, shuffling a deck, evaluating
an arithmetic expression. Ask in each case: can you recover the starting
state from the result? When the answer is no, ask what was lost and why.

The concept of aspect relegation is subtle. Frame it as: ``the work
happened earlier, and only its result is retained.'' Students who grasp
this will also grasp why expertise feels effortless and why it takes time
to develop---the effort is not reduced but relocated in time. This is the
single most important idea in Part~III. Take extra time here if needed.
\end{teachernote}

\begin{universitybridge}
Thermodynamics and the second law. Irreversibility and entropy
production. Partial evaluation and program specialisation. Computational
complexity and the limits of reversible computation (Bennett, Landauer).
\end{universitybridge}

\section*{Further Reading}

\begin{readinglist}
Charles H.~Bennett, \textit{Logical Reversibility of Computation}
(1973) --- connects computation and thermodynamics; short and precise.

Michael Sipser, \textit{Introduction to the Theory of Computation}
--- standard university CS text; includes discussion of time and
complexity.

Stephen Wolfram, \textit{A New Kind of Science} (selected sections)
--- explores computation unfolding over time with many visual examples.
\end{readinglist}
% ════════════════════════════════════════════════════════════════════════

\chapter{Local Coherence and Global Consistency}

\begin{historicalpreamble}
As programs became more complex, the nature of errors changed. Early
errors were often immediate and visible: a machine would halt, a number
would be obviously wrong. Later systems introduced subtle failures that
appeared correct under most conditions but failed under specific ones.
Debugging emerged as a discipline focused not only on identifying
mistakes, but on understanding why systems that seem correct can still
fail.

Formal verification methods were developed to prove that programs satisfy
all required conditions, not just the ones that happen to be tested.
These developments showed that correctness is not binary. A system can
appear correct while still violating constraints that are not immediately
visible. Local correctness---passing every individual test---does not
guarantee global consistency. This is one of the central structural facts
that formal verification was built to address, and it is the same fact
this chapter formalises.
\end{historicalpreamble}

\begin{historicalfigures}
Tony Hoare\index[topics]{Hoare, Tony} (1934--) developed Hoare logic, a formal system for reasoning
about program correctness. Donald Knuth\index[topics]{Knuth, Donald} (1938--) emphasised rigorous
analysis of algorithms and the importance of proving, not merely
assuming, correctness. Edsger Dijkstra advocated that programs should be
constructed in a way that makes correctness provable by design.
\end{historicalfigures}

\begin{paradigmtimeline}
1960s: debugging as a systematic discipline. \quad
1969: Hoare logic published. \quad
1970s: formal verification methods developed. \quad
1990s--2000s: static analysis tools and automated theorem provers.
\end{paradigmtimeline}

% ── Chapter Figure ────────────────────────────────────────────────────
\begin{figure}[h]
\centering
\begin{tikzpicture}[scale=0.97]
  % Sheet U1
  \fill[blue!18,cfigfill]
    (0,0) -- (3,0) -- (4,1) -- (1,1) -- cycle;
  \draw[cfig,blue!55!black]
    (0,0) -- (3,0) -- (4,1) -- (1,1) -- cycle;
  % Sheet U2
  \fill[green!18,cfigfill]
    (2,0.6) -- (5,0.6) -- (6,1.6) -- (3,1.6) -- cycle;
  \draw[cfig,green!50!black]
    (2,0.6) -- (5,0.6) -- (6,1.6) -- (3,1.6) -- cycle;
  % Sheet U3 (misaligned)
  \fill[red!18,cfigfill]
    (3.4,1.3) -- (6.4,1.5) -- (7.3,2.4) -- (4.3,2.2) -- cycle;
  \draw[cfig,red!55!black]
    (3.4,1.3) -- (6.4,1.5) -- (7.3,2.4) -- (4.3,2.2) -- cycle;
  % Mismatch seam
  \draw[very thick,red!70,dashed] (3.7,1.4) -- (4.85,1.52);
  % Section labels
  \node[blue!65!black]  at (1.5,-0.45) {$U_1$};
  \node[green!55!black] at (4.3, 0.2)  {$U_2$};
  \node[red!65!black]   at (6.0, 2.85) {$U_3$};
  \node[red!70!black]   at (5.15, 1.1) {gluing fails};
  % Overlap indicator
  \draw[<->,thin,gray!60] (3.65,0.78) -- (3.65,1.35);
  \node[gray!70!black] at (2.95,1.0) {overlap $U_1\!\cap\! U_2$};
\end{tikzpicture}
\caption{Local sections $s_1, s_2, s_3$ each satisfy constraints on their own region.
At the overlap of $U_2$ and $U_3$ the sections disagree: gluing fails.
Local correctness does not guarantee global consistency.}
\end{figure}

\section{The Local-Global Distinction}

\begin{definition}[Local Consistency]\index[vocab]{local consistency}
A configuration is \textbf{locally consistent} within a region $U \subset X$
if it satisfies all constraints $c \in \C$ restricted to $U$.
\end{definition}

\begin{definition}[Global Consistency]\index[vocab]{global consistency}
A configuration is \textbf{globally consistent} if it satisfies all
constraints $c \in \C$ on the entire domain $X$.
\end{definition}

Local consistency does not imply global consistency. This is one of the
central facts about structural thinking, and its failure modes are common
and important.

\section{Fragmentation}

\begin{definition}[Fragmented Structure]\index[vocab]{fragmentation}
A structure is \textbf{fragmented} if it is locally consistent on each
of several regions but fails global consistency when the regions are
combined.
\end{definition}

\begin{example}
The Penrose staircase\index[topics]{Penrose staircase} appears locally consistent at every corner: each step
connects to the next. But globally, it is impossible---following the stairs
returns you to the starting point after ascending. The global section fails.
\end{example}

Fragmentation is a common failure in complex theories. Each piece makes
local sense. The contradiction only appears when pieces are combined.

\section{Overcompression}

\begin{definition}[Overcompressive Abstraction]\index[vocab]{overcompression}
An abstraction $\alpha : X \to Y$ is \textbf{overcompressive} if it
identifies configurations that are distinguishable under $\C$:
\[
  \exists\, x_1, x_2 \in X : \alpha(x_1) = \alpha(x_2)
  \text{ but } x_1 \not\sim_\C x_2.
\]
\end{definition}

Overcompression discards relevant distinctions. The resulting model appears
simpler and may even appear internally consistent, because the distinctions
that would reveal the inconsistency have been erased.

\begin{proposition}
An overcompressive abstraction generically increases the obstruction to
global consistency, because the identified configurations impose incompatible
constraints at their shared boundaries.
\end{proposition}

\section{Hallucination as False Global Section}

\begin{definition}[Hallucinated Structure]\index[vocab]{hallucination}
A \textbf{hallucinated structure} is a configuration that appears globally
consistent under a restricted projection but does not correspond to any
element of $\A(\C)$ for the full constraint set.
\end{definition}

Hallucination is overcompression producing a false sense of closure. The
constraints that would reveal the impossibility have been removed from
consideration.

\begin{remark}
This definition applies equally to human belief formation and to artificial
systems. A language model that produces a confident but false statement is
exhibiting hallucination in precisely this formal sense: the local consistency
within its training distribution does not extend to global consistency with
the full constraint set of reality.
\end{remark}

% ── Formal Correspondence ─────────────────────────────────────────────
\section*{Formal Correspondence}
\begin{formalcorrespondence}
Let $\{U_i\}$ be overlapping domains of definition with local sections
$s_i$ valid on each $U_i$. Global coherence requires compatibility on
overlaps:
\[
  s_i\big|_{U_i \cap U_j} = s_j\big|_{U_i \cap U_j}.
\]
If this holds, the sections glue into a global section $s$. Failure of
this condition corresponds to inconsistency: local correctness without
global realisability. A system of constraints $\mathcal{C}$ is consistent
if $\Gamma_{\mathcal{C}} \neq \emptyset$ and inconsistent if
$\Gamma_{\mathcal{C}} = \emptyset$. This is the sheaf-theoretic restatement
of the definitions above.
\end{formalcorrespondence}

% ── Computational Aside ───────────────────────────────────────────────
\section*{Computational Perspective}
\begin{computationalaside}
The most difficult bugs are not immediate failures but near-successes. A
system appears correct under limited testing but fails under more complete
conditions. Ideas fail in the same way. Structures that satisfy some
constraints but not all can appear stable until examined more fully.
Failure is often the result of incomplete constraint resolution rather
than obvious contradiction.
\end{computationalaside}

% ── Worked Example ────────────────────────────────────────────────────
\section*{Worked Example}
Consider three local constraints: $C_1: x = y$, $C_2: y = z$,
$C_3: x \neq z$. Each pair appears locally consistent. But $C_1$ and
$C_2$ together imply $x = y = z$, which directly contradicts $C_3$.
Thus $\Gamma_{\mathcal{C}} = \emptyset$. The system is locally coherent
but globally inconsistent. Local satisfaction of constraints does not
guarantee global realisability.

\begin{exercise}
Imagine two people editing the same document without communicating.
Describe how inconsistencies could arise even if both make reasonable
changes. What would be required to ensure a consistent final result?
\end{exercise}

% ── Spherepop Comparison ──────────────────────────────────────────────
\section*{Spherepop Comparison}
\begin{spherepopcomparison}
A Spherepop system can look valid locally while failing globally when
different reduced regions impose incompatible outer boundaries.

\textbf{BNF fragment.}
\begin{verbatim}
<tile>    ::= "tile(" <region> ")"
<overlap> ::= <tile> "&" <tile>
<global>  ::= "glue(" <tile-list> ")"
<failure> ::= "mismatch(" <overlap> ")"
\end{verbatim}

Fragmentation occurs when tiles exist but do not glue. Hallucination
occurs when a seeming whole has no admissible underlying global region.
These correspond directly to the sheaf gluing condition stated above.
\end{spherepopcomparison}

\begin{umldiagram}
  \node[box] (l1) {Local\\Region $A$};
  \node[box, below=1.4cm of l1] (l2) {Local\\Region $B$};
  \node[box, right=3.2cm of $(l1)!0.5!(l2)$] (gl) {Global\\Glue Attempt};
  \draw[arrow] (l1) -- (gl);
  \draw[arrow] (l2) -- (gl);
\end{umldiagram}

% ── Teacher Notes ─────────────────────────────────────────────────────
\section*{Teacher Notes}

\begin{teachernote}
Students must understand that local correctness is not sufficient for
global coherence, and that this is a structural fact, not a matter of
carelessness. The Penrose staircase example in the chapter body is
useful precisely because it makes the failure vivid and unambiguous.
Use it early.

The most common difficulty is students accepting partial solutions as
complete. Design exercises that require students to combine locally
valid pieces and check whether the combination is globally consistent.
Give systems that almost work: each piece checks out, but the pieces
disagree on their shared boundary.

Watch for students who fail to check overlap conditions. Ask explicitly:
does your solution work on the \emph{intersection}? This habit of
checking overlaps---not just individual pieces---is the central skill
this chapter builds, and it generalises directly to proof verification,
software testing, and scientific theory-building.
\end{teachernote}

\begin{universitybridge}
Sheaf theory and \v{C}ech cohomology. Distributed systems and
consistency models (CAP theorem, eventual consistency). Constraint
propagation. Formal verification and model checking.
\end{universitybridge}

\section*{Further Reading}

\begin{readinglist}
Leslie Lamport, \textit{Time, Clocks, and the Ordering of Events in a
Distributed System} (1978) --- a classic paper readable at advanced
high school level with guidance.

Eric Brewer, overview articles on the CAP theorem --- accessible
survey of the trade-offs in distributed consistency.

Glen E.~Bredon, \textit{Sheaf Theory} (introductory ideas) --- the
formal mathematical treatment at university level.
\end{readinglist}

\chapter{Cognitive Dissonance as Obstruction}

\begin{historicalpreamble}
Grace Hopper's team built the first compiler in the early 1950s, and with
it came a new insight: computation could be transformed before it ran.
A compiler does not execute a program; it translates it into a form that
a machine can execute more efficiently, performing much of the work in
advance. Optimisations remove redundant steps, reorganise computations,
and produce a representation that runs quickly.

The result is a system that appears to ``know what to do'' immediately,
even though the work was carried out in a prior stage. This distinction
between interpreted and compiled execution revealed that speed often comes
from prior transformation rather than inherent simplicity---and it is
exactly the distinction between explicit reasoning and intuition that
this chapter formalises as an obstruction problem.
\end{historicalpreamble}

\begin{historicalfigures}
Grace Hopper\index[topics]{Hopper, Grace} (1906--1992) developed the first compiler and coined the
term ``debugging.'' John McCarthy\index[topics]{McCarthy, John} (1927--2011) advanced symbolic
computation and introduced the idea of programs that reason about their
own structure. Frances Allen\index[topics]{Allen, Frances} (1932--2020) pioneered compiler optimisation
techniques that became foundational to modern computing.
\end{historicalfigures}

\begin{paradigmtimeline}
1950s: early compilers (Hopper). \quad
1970s: optimisation techniques formalised. \quad
1990s: just-in-time compilation. \quad
2000s: adaptive and profile-guided optimising systems.
\end{paradigmtimeline}

% ── Chapter Figure ────────────────────────────────────────────────────
\begin{figure}[h]
\centering
\begin{tikzpicture}[scale=0.97]
  % Energy landscape curve
  \draw[very thick,gray!70]
    (0,2.5)
    .. controls (1,1.0)  and (2,1.0)  .. (3,2.4)
    .. controls (3.8,3.2) and (4.6,3.2) .. (5.5,2.4)
    .. controls (6.4,1.0) and (7.4,1.0) .. (8.5,2.5);
  % Attractor points
  \fill[blue!65] (1.45,1.42) circle (2.5pt);
  \fill[blue!65] (6.95,1.42) circle (2.5pt);
  \node at (1.45,0.95) {$A$};
  \node at (6.95,0.95) {$B$};
  % Energetic barrier label
  \node at (4.25,3.55) {energetic barrier};
  % Curved arrow (energetic path — possible but costly)
  \draw[->,thick,blue!50]
    (1.85,1.75) .. controls (3.0,3.1) and (5.5,3.1) .. (6.65,1.75);
  \node[blue!65!black] at (4.25,3.1) {};
  % Topological cut (cannot cross)
  \draw[very thick,red!65] (4.25,0.2) -- (4.25,-1.3);
  \node[red!65!black] at (5.9,-0.65) {topological separation};
  % Axis
  \draw[->] (-0.3,0) -- (8.9,0) node[right]{configuration};
  \draw[->] (0,-0.2) -- (0,4.0) node[above]{$E$};
\end{tikzpicture}
\caption{States $A$ and $B$ occupy separate minima. They are energetically separated if
a continuous path through the landscape connects them. They are topologically separated
if no such path exists: no amount of gradual change can bridge the gap.}
\end{figure}

\section{A Precise Definition}

\begin{definition}[Gluing Failure]\index[vocab]{gluing failure}
Given two locally consistent structures $\psi_1$ on region $U_1$ and
$\psi_2$ on region $U_2$ with $U_1 \cap U_2 \neq \emptyset$, a
\textbf{gluing failure} occurs when $\psi_1$ and $\psi_2$ disagree on
$U_1 \cap U_2$---that is, they cannot be combined into a single globally
consistent structure.
\end{definition}

Cognitive dissonance is the experiential signal of a gluing failure. Two
beliefs, each locally coherent, cannot be reconciled at their overlap.

\section{Energetic and Topological Separation}

\begin{definition}[Energetic Separation]\index[vocab]{energetic separation}
Two locally consistent configurations are \textbf{energetically separated}
if they can be connected by a continuous path of intermediate configurations,
each locally consistent, at some cost. Transition between them is possible
by gradual change.
\end{definition}

\begin{definition}[Topological Separation]\index[vocab]{topological separation}
Two configurations are \textbf{topologically separated} if they lie in
distinct connected components of the space of globally consistent
structures. No continuous path of locally consistent configurations
connects them.
\end{definition}

This distinction has a direct experiential analog:

\begin{itemize}
  \item Energetically separated beliefs can be reconciled by gradual
        accumulation of evidence. The transition is continuous.
  \item Topologically separated beliefs require a discontinuous
        restructuring---a paradigm shift, a gestalt change. No amount of
        gradual evidence accumulation reaches across the topological barrier.
\end{itemize}

\begin{exercise}
Think of a belief you hold. Is it energetically or topologically separated
from its opposite? What would it take to change it: gradual evidence, or a
sudden restructuring of your entire framework?
\end{exercise}

% ── Formal Correspondence ─────────────────────────────────────────────
\section*{Formal Correspondence}
\begin{formalcorrespondence}
Let $T$ be a sequence of constraint applications. Explicit reasoning
corresponds to stepwise evaluation $x_{n+1} = T(x_n)$. Intuition
corresponds to a precomputed mapping $x \mapsto x^*$ where $x^*$ is the
result of many implicit iterations. Thus intuition is a compressed
representation of iterative computation. Cognitive dissonance is the
obstruction that arises when two such precomputed mappings produce
incompatible results on their shared domain---a gluing failure in the
space of compressed representations.
\end{formalcorrespondence}

% ── Computational Aside ───────────────────────────────────────────────
\section*{Computational Perspective}
\begin{computationalaside}
In distributed systems, two components may each function correctly in
isolation, yet the system as a whole fails if their states do not align.
Local correctness does not ensure global consistency. Understanding
behaves the same way: a collection of correct parts is not sufficient;
they must agree on their overlaps. Coherence is a condition on the whole,
not just the parts.
\end{computationalaside}

% ── Worked Example ────────────────────────────────────────────────────
\section*{Worked Example}
Consider two constraint systems: $\mathcal{C}_A: x > 5$ and
$\mathcal{C}_B: x < 3$. Individually, each admits solutions:
$\Gamma_A = \{6,7,8,\dots\}$ and $\Gamma_B = \{1,2\}$. Together,
$\Gamma_{A \cup B} = \emptyset$. The incompatibility cannot be resolved
by refinement alone. This is an obstruction: two valid structures that
cannot be merged into a single consistent configuration. Whether the
separation is energetic or topological determines what resolution, if
any, is possible.

% ── Spherepop Comparison ──────────────────────────────────────────────
\section*{Spherepop Comparison}
\begin{spherepopcomparison}
In Spherepop, obstruction appears when two or more region-structures
cannot be reduced into a single coherent whole without tearing the
grammar of composition.

\textbf{BNF fragment.}
\begin{verbatim}
<pathA>       ::= <region> "=>" ... "=>" <region>
<pathB>       ::= <region> "=>" ... "=>" <region>
<obstruction> ::= "obs(" <pathA> "," <pathB> ")"
\end{verbatim}

Topological separation means there is no admissible continuous reduction
path between the two region-families. This is the grammar-level
expression of gluing failure---the same condition stated in the
definitions above.
\end{spherepopcomparison}

\begin{umldiagram}
  \node[box] (s1) {Stable\\Region 1};
  \node[box, right=4cm of s1] (s2) {Stable\\Region 2};
  \node[box, below=1.8cm of $(s1)!0.5!(s2)$] (obs) {Obstruction /\\Separation};
  \draw[arrow] (s1) -- (obs);
  \draw[arrow] (s2) -- (obs);
\end{umldiagram}

% ── Teacher Notes ─────────────────────────────────────────────────────
\section*{Teacher Notes}

\begin{teachernote}
Students should recognise when systems cannot be reconciled, and accept
that recognition as a result, not a failure. The most common difficulty
is a refusal to acknowledge genuine obstruction: students try to force
incompatible ideas together by weakening one or reinterpreting both. It
is important to establish clearly that some constraint systems have no
solution, and that demonstrating this is a legitimate and useful
conclusion.

The distinction between energetic and topological separation is worth
spending time on. Energetic separation means that gradual change can
bridge the gap; topological separation means it cannot. Ask students to
think of a belief they would change with enough evidence, and a belief
they would not change regardless of evidence. The second type is
topologically separated from its negation. This is not a comment on the
belief's truth, but on the structure of the space it inhabits.

Watch for students who believe more effort resolves contradiction.
The most important thing this chapter teaches is the skill of recognising
genuine incompatibility early, rather than investing effort in resolution
that is structurally impossible.
\end{teachernote}

\begin{universitybridge}
Topological obstruction and cohomology. Unsatisfiable constraint systems
and UNSAT certificates. G\"odel's incompleteness theorems (conceptually).
Logical inconsistency and proof by contradiction.
\end{universitybridge}

\section*{Further Reading}

\begin{readinglist}
Raymond Smullyan, \textit{What Is the Name of This Book?} ---
logical paradoxes and contradictions presented accessibly.

Edsger Dijkstra, \textit{A Discipline of Programming} --- rigorous
reasoning about program correctness and the structure of logical
argument.

Allen Hatcher, \textit{Algebraic Topology} (conceptual introduction
only) --- introduces obstruction in a geometric setting at university
level.
\end{readinglist}
% ════════════════════════════════════════════════════════════════════════

\chapter{Functorial Correspondence}

\begin{historicalpreamble}
As data grew in scale, storing and transmitting it efficiently became a
central problem. Compression algorithms were developed to reduce the size
of data while preserving its essential structure. Some methods, such as
lossless compression, preserve all information. Others, such as lossy
compression, discard details unlikely to be noticed, retaining only what
is most important.

Claude Shannon's 1948 paper founding information theory showed that the
compressibility of a message is determined not by its content but by its
structure---specifically, by the statistical redundancy of the symbols it
uses. Encoding schemes translate complex structures into simpler
representations that can be stored, transmitted, or processed more
efficiently. These developments showed that representation is always a
reduction of underlying complexity, and that preserving structure is more
important than preserving surface form. This is the same insight that
functorial correspondence formalises across mathematical domains.
\end{historicalpreamble}

\begin{historicalfigures}
Claude Shannon\index[topics]{Shannon, Claude} (1916--2001) founded information theory with his 1948
paper ``A Mathematical Theory of Communication.'' David Huffman
(1925--1999) developed the optimal prefix-free compression code bearing
his name. Abraham Lempel\index[topics]{Lempel, Abraham} (1936--) and Jacob Ziv\index[topics]{Ziv, Jacob} (1931--) created the
LZ family of compression algorithms, which underlie most lossless
compression used today.
\end{historicalfigures}

\begin{paradigmtimeline}
1948: Shannon's information theory paper. \quad
1950s: coding theory and error-correcting codes. \quad
1970s: Lempel--Ziv compression algorithms. \quad
2000s: multimedia encoding, streaming, and perceptual compression.
\end{paradigmtimeline}

% ── Chapter Figure ────────────────────────────────────────────────────
\begin{figure}[h]
\centering
\begin{tikzpicture}[scale=0.93]
  % Left cube
  \draw[cfig] (0,0) -- (2,0) -- (2,2) -- (0,2) -- cycle;
  \draw[cfig] (0.7,0.7) -- (2.7,0.7) -- (2.7,2.7) -- (0.7,2.7) -- cycle;
  \draw[cfig] (0,0)--(0.7,0.7); \draw[cfig] (2,0)--(2.7,0.7);
  \draw[cfig] (2,2)--(2.7,2.7); \draw[cfig] (0,2)--(0.7,2.7);
  % Vertex labels (left)
  \node[cfignode] at (0,0)     {}; \node at (-0.2,-0.2) {$a$};
  \node[cfignode] at (2,0)     {}; \node at ( 2.2,-0.2) {$b$};
  \node[cfignode] at (2,2)     {}; \node at ( 2.2, 2.2) {$c$};
  \node[cfignode] at (2.7,2.7) {}; \node at ( 2.9, 2.9) {$d$};
  \node at (1.35,-0.6) {domain $\mathcal{C}$};
  % Right cube
  \draw[cfig] (6,0) -- (8,0) -- (8,2) -- (6,2) -- cycle;
  \draw[cfig] (6.7,0.7) -- (8.7,0.7) -- (8.7,2.7) -- (6.7,2.7) -- cycle;
  \draw[cfig] (6,0)--(6.7,0.7); \draw[cfig] (8,0)--(8.7,0.7);
  \draw[cfig] (8,2)--(8.7,2.7); \draw[cfig] (6,2)--(6.7,2.7);
  % Vertex labels (right)
  \node[cfignode] at (6,0)     {}; \node at (5.8,-0.2) {$F(a)$};
  \node[cfignode] at (8,0)     {}; \node at (8.25,-0.2){$F(b)$};
  \node[cfignode] at (8,2)     {}; \node at (8.25, 2.2){$F(c)$};
  \node[cfignode] at (8.7,2.7) {}; \node at (8.95, 2.9){$F(d)$};
  \node at (7.35,-0.6) {domain $\mathcal{D}$};
  % Structure-preserving arrows
  \draw[->,thick,blue!60] (2.7,2.7) to[bend left=12] (6.0,2.7);
  \draw[->,thick,blue!60] (2.0,2.0) to[bend left=8]  (8.0,2.0);
  \draw[->,thick,blue!60] (0.7,0.7) to[bend right=8] (6.7,0.7);
  \node[blue!70!black] at (4.35,3.3) {$F$};
\end{tikzpicture}
\caption{A structure-preserving map $F$ carries both vertices and their compositional
relations from domain $\mathcal{C}$ into domain $\mathcal{D}$.
What changes is the interpretation of objects; what is preserved is the grammar of their relationships.}
\end{figure}

\section{Structure-Preserving Maps}

\begin{definition}[Structure-Preserving Map]\index[vocab]{structure-preserving map}\index[vocab]{functor}
A map $F : D_1 \to D_2$ between domains is \textbf{structure-preserving}
if it maps admissible configurations to admissible configurations and
preserves the compositional rules of the domain:
\[
  x \in \A(\C_1) \;\Rightarrow\; F(x) \in \A(\C_2),
\]
and if $x_1 \circ x_2$ is an admissible composition in $D_1$, then
$F(x_1) \circ F(x_2)$ is admissible in $D_2$.
\end{definition}

In the language of category theory, such maps are called functors. For our
purposes, the key point is simpler: a structure-preserving map translates
the constraint grammar of one domain into another without distorting it.

\section{What Is Preserved}

\begin{theorem}[Invariants Under Structure-Preserving Maps]\index[topics]{functorial correspondence, invariants}
A structure-preserving map $F : D_1 \to D_2$ preserves:
\begin{enumerate}[label=(\roman*)]
  \item Admissibility: elements of $\A(\C_1)$ map to $\A(\C_2)$.
  \item Obstruction: gluing failures in $D_1$ produce gluing failures
        in $D_2$.
  \item Attractor structure: stable configurations in $D_1$ map to
        stable configurations in $D_2$.
\end{enumerate}
\end{theorem}

\begin{remark}
What changes under $F$ is the interpretation of the objects. What does not
change is the grammar of their relationships. This is why the same
mathematical structures appear in physics, cognition, and economics: the
grammar is invariant; the objects are different.
\end{remark}

\section{Examples}

\begin{example}
The flow of heat through a solid and the flow of electrical current through
a conductor are related by a structure-preserving map. Temperature corresponds
to voltage, heat flux to current density, thermal conductivity to electrical
conductivity. The equations governing one can be translated directly into
the equations governing the other.
\end{example}

\begin{example}
The spread of information through a social network and the spread of a
contagion through a population share a structure-preserving relationship.
Connectivity, transmission rate, and recovery correspond across domains.
Insights from epidemiology directly inform information theory and vice versa.
\end{example}

% ── Formal Correspondence ─────────────────────────────────────────────
\section*{Formal Correspondence}
\begin{formalcorrespondence}
Let $X$ be a high-dimensional state space. A projection $\pi : X \to Y$
reduces dimensionality. Multiple states may map to the same observable:
$\pi(x_1) = \pi(x_2)$. The projection preserves invariants but loses
internal detail. A structure-preserving map $F : D_1 \to D_2$ does the
same across domains: it translates the constraint grammar of one domain
into another without distorting it. What changes under $F$ is the
interpretation of the objects; what does not change is the structure of
their relationships. This corresponds to coarse-graining in physical and
informational systems, and to the functor concept in category theory.
\end{formalcorrespondence}

% ── Computational Aside ───────────────────────────────────────────────
\section*{Computational Perspective}
\begin{computationalaside}
Compression reduces a representation while preserving structure necessary
for reconstruction or use. Projection behaves similarly: multiple
underlying configurations may be represented by a single observable
state. Understanding depends on selecting representations that preserve
relevant structure while discarding unnecessary detail. The same
principle underlies both Shannon's compression theory and the
category-theoretic notion of a functor.
\end{computationalaside}

% ── Worked Example ────────────────────────────────────────────────────
\section*{Worked Example}
Consider two domains. Numbers: $1, 2, 3$. Shapes: $\triangle, \square,
\bigcirc$. Define a mapping $F: 1 \mapsto \triangle,\ 2 \mapsto \square,\
3 \mapsto \bigcirc$. If addition in numbers corresponds to composition of
shapes---$1 + 2 = 3$ corresponds to $\triangle \circ \square = \bigcirc$---
then structure is preserved across domains. What the map preserves is not
the objects themselves but the relationships between them.

\begin{exercise}
Find two domains in your own experience where the same structural
relationship holds. Describe the mapping explicitly. What is preserved,
and what changes?
\end{exercise}

% ── Spherepop Comparison ──────────────────────────────────────────────
\section*{Spherepop Comparison}
\begin{spherepopcomparison}
Functorial correspondence means that Spherepop structures may be mapped
into another domain while preserving the grammar of composition.

\textbf{BNF fragment.}
\begin{verbatim}
<object>   ::= <region>
<morphism> ::= <object> "->" <object>
<functor>  ::= "F(" <object> ")" | "F(" <morphism> ")"
\end{verbatim}

If a merge or collapse relation exists in Spherepop, a
structure-preserving map sends it to a valid relation in the target
domain. The map preserves admissibility, obstruction, and attractor
structure---the same invariants listed in the theorem above.
\end{spherepopcomparison}

\begin{umldiagram}
  \node[box] (src) {Spherepop\\Objects \& Morphisms};
  \node[box, right=3.2cm of src] (tgt) {Target Domain\\Objects \& Morphisms};
  \draw[arrow] (src) -- node[above]{\small Functor $F$} (tgt);
\end{umldiagram}

% ── Teacher Notes ─────────────────────────────────────────────────────
\section*{Teacher Notes}

\begin{teachernote}
Students should come to see analogy not as a rhetorical device but as a
structure-preserving mapping that can be verified. The most common
difficulty is superficial analogy: students map surface features---names,
appearances, casual resemblances---rather than structural relationships.
A well-formed analogy maps morphisms, not just objects.

Design exercises that require students to state explicitly what the
mapping preserves. Ask: if I perform this operation on one side, what
must the corresponding operation be on the other? If they cannot answer,
the analogy is not yet structural.

Watch for students who map objects when they should be mapping
relationships. The key question is always: is the composition rule
preserved? If $f$ followed by $g$ produces $h$ in one domain, does
$F(f)$ followed by $F(g)$ produce $F(h)$ in the other? That is the
criterion for a functorial correspondence, and it is the test that
separates genuine structural analogy from coincidental similarity.
\end{teachernote}

\begin{universitybridge}
Category theory and functors. Representation theory. Formal analogies
in physics (e.g., heat and electricity). Cross-domain modelling in
applied mathematics.
\end{universitybridge}

\section*{Further Reading}

\begin{readinglist}
F.~William Lawvere and Stephen Schanuel, \textit{Conceptual
Mathematics} --- the most accessible introduction to category theory;
suitable for advanced high school students.

Saunders Mac Lane, \textit{Categories for the Working Mathematician}
--- the standard university-level reference for functors and natural
transformations.

David I.~Spivak, \textit{Category Theory for the Sciences} ---
applications of categorical structure across scientific domains.
\end{readinglist}

\chapter{Intelligence as Constraint Maintenance}

\begin{historicalpreamble}
Object-oriented programming emerged as a way to manage complexity by
organising software around interacting components. Each object combines
data and behaviour, encapsulating its internal state while exposing
methods for interaction. Systems built in this way emphasise
relationships: objects communicate, collaborate, and depend on one
another to produce overall behaviour.

This approach shifted the focus from individual operations to the
structure of interactions. The behaviour of the system is no longer
found in a single place but emerges from the coordinated activity of
many parts. Understanding such a system requires understanding how
components relate, not just what they contain. This is the same shift
that this chapter makes in the theory of intelligence: from intelligence
as internal optimisation to intelligence as the maintenance of
constraint-consistent structure across a network of relationships.
\end{historicalpreamble}

\begin{historicalfigures}
Alan Kay (1940--) pioneered object-oriented programming and the idea
that meaning lies in interfaces between components rather than in their
internal construction. Niklaus Wirth\index[topics]{Wirth, Niklaus} (1934--2024) designed structured
languages that made compositional reasoning tractable. Bjarne Stroustrup
(1950--) developed C++, bringing object-oriented design into systems
programming and shaping how large-scale software is structured today.
\end{historicalfigures}

\begin{paradigmtimeline}
1970s: object-oriented programming formalised. \quad
1983: C++ developed. \quad
1995: Java and platform-independent component systems. \quad
2000s: component-based and service-oriented architectures.
\end{paradigmtimeline}

% ── Chapter Figure ────────────────────────────────────────────────────
\begin{figure}[h]
\centering
\begin{tikzpicture}[scale=0.97]
  % Corridor walls
  \draw[cfig,gray!60]
    (0,0.55) .. controls (2,1.05) and (4,1.45) .. (8.5,1.65);
  \draw[cfig,gray!60]
    (0,-0.55) .. controls (2,-0.2) and (4,0.1) .. (8.5,0.38);
  % Corridor fill
  \fill[blue!8]
    (0,0.55) .. controls (2,1.05) and (4,1.45) .. (8.5,1.65)
    -- (8.5,0.38) .. controls (4,0.1) and (2,-0.2) .. (0,-0.55) -- cycle;
  % Admissible trajectory
  \draw[very thick,blue!65]
    (0,0) .. controls (2,0.22) and (4,0.82) .. (8.5,1.02);
  \fill[blue!65] (8.5,1.02) circle (2pt);
  \node[blue!70!black] at (7.2,0.6) {admissible trajectory};
  % Unconstrained optimum (dashed, exits corridor)
  \draw[very thick,dashed,red!60]
    (4,0.82) .. controls (5,2.9) and (6.2,3.1) .. (7.8,2.3);
  \fill[red!60] (7.8,2.3) circle (2pt);
  \node[red!65!black] at (8.25,2.45) {unconstrained optimum};
  % Label
  \node at (4.25,-0.95) {constraint corridor};
  % Axis
  \draw[->] (-0.3,-1.05) -- (9.0,-1.05) node[right]{$t$};
\end{tikzpicture}
\caption{The admissible trajectory stays within the constraint corridor at all times.
The unconstrained optimum lies outside the corridor: pursuing it would violate the
conditions required for the system's own continuation.}
\end{figure}

\section{The Optimization Model and Its Failure}

The dominant model of intelligence---in both cognitive science and artificial
intelligence---treats intelligent systems as optimizers: agents that maximize
an objective function.

This model is useful but incomplete. It fails in predictable ways when:
\begin{itemize}
  \item The objective function does not capture all relevant constraints.
  \item The system's actions alter the constraint environment itself.
  \item Optimization of the objective is incompatible with the persistence
        of the system.
\end{itemize}

\begin{principle}[Intelligence as Constraint Maintenance]\index[topics]{intelligence, as constraint maintenance}
An intelligent system is one that can maintain constraint-consistent
structure across time, acquiring new constraints, revising inconsistent
ones, and preserving globally coherent configurations as the environment
changes.
\end{principle}

\section{Non-Admissibility of Self-Invalidating Trajectories}

\begin{definition}[Constraint-Preserving Evolution]\index[vocab]{constraint-preserving evolution}
A trajectory $\gamma(t)$ is \textbf{constraint-preserving} if
$\gamma(t) \in \A(\C)$ for all $t$.
\end{definition}

\begin{definition}[Substrate Violation]\index[vocab]{substrate violation}
A trajectory \textbf{violates its substrate} if it alters $\C$ such that
future states are no longer admissible under the constraints required for
the system's own persistence.
\end{definition}

\begin{theorem}[Non-Admissibility of Self-Invalidating Trajectories]\index[topics]{self-invalidating trajectories}
A trajectory that leads to a state where the constraint set $\C$ required
for its continuation is no longer satisfied is not in $\A(\C)$.
\end{theorem}

\begin{proof}
Admissibility requires that each continuation of $\gamma$ satisfies $\C$.
If $\gamma(t^*)$ violates $\C$, no extension of $\gamma$ exists within
$\A(\C)$. Therefore $\gamma \notin \A(\C)$.
\end{proof}

\begin{corollary}
Any trajectory that eliminates the conditions required for its own operation
is not a stable solution of the system. ``Doom scenarios'' in which an
autonomous system converts all available resources into forms incompatible
with its own persistence are non-admissible by this theorem.
\end{corollary}

\section{Bounded Self-Improvement}

\begin{theorem}[Constraint-Bounded Self-Modification]\index[topics]{bounded self-improvement}
Let $\mathcal{R}$ be a self-modification operator that is constraint-preserving.
Then the sequence $\{\mathcal{R}^n\}$ of iterated self-modifications remains
within $\A(\C)$.
\end{theorem}

\begin{proof}
Each application of $\mathcal{R}$ maps $\A(\C)$ into $\A(\C)$ by assumption.
The composition of constraint-preserving maps is constraint-preserving.
Therefore no finite iteration escapes $\A(\C)$.
\end{proof}

\begin{corollary}
Unbounded self-improvement is not possible for a constraint-preserving
system. Each increment of modification must satisfy the same constraints
as the system being modified.
\end{corollary}

\begin{openproblem}
Characterize the conditions under which real-world systems can be
approximated as constraint-preserving. What perturbations to the
constraint environment are tolerable before the approximation breaks down?
\end{openproblem}

% ── Formal Correspondence ─────────────────────────────────────────────
\section*{Formal Correspondence}
\begin{formalcorrespondence}
A system can be modelled as objects and morphisms:
\[
  A \xrightarrow{f} B \xrightarrow{g} C, \qquad g \circ f : A \to C.
\]
Structure is defined not by the internal details of objects but by
relationships preserved under composition. Intelligence as constraint
maintenance corresponds to a trajectory $\gamma(t) \in \A(\C)$ for all
$t$---a path through object-space that respects all morphisms. A
self-invalidating trajectory violates this condition: it destroys the
relationships required for its own continuation. The theorem of
non-admissibility states that such trajectories are not elements of
$\A(\C)$.
\end{formalcorrespondence}

% ── Computational Aside ───────────────────────────────────────────────
\section*{Computational Perspective}
\begin{computationalaside}
In object-oriented systems, components are defined by how they interact
rather than by their internal structure. The behaviour of the system
emerges from relationships between parts. Knowledge follows the same
pattern: meaning is determined not by internal composition alone, but by
how structures participate in a network of interactions. Intelligence
is the capacity to maintain that network under changing conditions.
\end{computationalaside}

% ── Worked Example ────────────────────────────────────────────────────
\section*{Worked Example}
Consider a system minimising error $E(x) = (x - 5)^2$. The unconstrained
minimum is $x = 5$. Now introduce a conflicting constraint $x \geq 7$.
The new admissible solution is $x^* = 7$: the system maintains
constraints by selecting the best admissible approximation rather than
the unconstrained optimum. Intelligence consists in maintaining
consistency under competing conditions, not in ignoring constraints to
maximise a single objective.

\begin{exercise}
Break a complex system---a school, a game, a supply chain---into
interacting parts. Describe how the behaviour of the whole depends on
relationships between parts rather than on the internal structure of any
single part.
\end{exercise}

% ── Spherepop Comparison ──────────────────────────────────────────────
\section*{Spherepop Comparison}
\begin{spherepopcomparison}
Spherepop makes intelligence look less like unconstrained optimisation
and more like the maintenance of an admissible nested structure.

\textbf{BNF fragment.}
\begin{verbatim}
<agent>      ::= "agent(" <region> ")"
<trajectory> ::= <agent> "->" <agent> "->" <agent>
<invalid>    ::= "break(" <trajectory> ")"
\end{verbatim}

A self-invalidating trajectory is one in which a Spherepop structure
destroys the outer conditions required for further composition. Such a
path is not a higher success but a structural failure---non-admissible
by the same criterion as the theorem above.
\end{spherepopcomparison}

\begin{umldiagram}
  \node[box] (ag)  {Agent as\\Trajectory};
  \node[box, right=2.4cm of ag]  (maint) {Constraint\\Maintenance};
  \node[box, right=2.4cm of maint] (stab) {Admissible\\Intelligence};
  \draw[arrow] (ag) -- (maint);
  \draw[arrow] (maint) -- (stab);
\end{umldiagram}

% ── Teacher Notes ─────────────────────────────────────────────────────
\section*{Teacher Notes}

\begin{teachernote}
Students should understand intelligence as the maintenance of
constraint-consistent structure over time, not as the maximisation of a
single objective. The most common difficulty is equating intelligence
with optimisation: students accept that an agent should maximise
something, and then ask which thing. Resist this framing. Begin instead
by asking: what conditions must a system satisfy to continue existing
and operating? Those conditions are the constraints. Intelligence is
what maintains them.

Introduce conflicting constraints explicitly. Give scenarios where
maximising one objective destroys the conditions required to pursue it.
The corollary about self-invalidating trajectories should feel
inevitable once students have worked through such a scenario themselves.

Watch for students choosing solutions that are optimal under a single
objective but inadmissible under the full constraint set. Ask them to
check each candidate against all constraints before accepting it. The
habit of checking all constraints, not just the most salient one, is
the central skill of this chapter.
\end{teachernote}

\begin{universitybridge}
Constrained optimisation and Lagrange multipliers. Control theory and
stability. AI alignment and value specification. Dynamical systems
stability and Lyapunov functions.
\end{universitybridge}

\section*{Further Reading}

\begin{readinglist}
Stuart Russell and Peter Norvig, \textit{Artificial Intelligence: A
Modern Approach} --- the standard AI textbook; covers constraint
satisfaction, optimisation, and decision systems.

Karl Friston, introductory papers on the free-energy principle ---
describes intelligence as constraint maintenance under a variational
formulation; accessible survey articles available.

Richard Sutton and Andrew Barto, \textit{Reinforcement Learning: An
Introduction} --- optimisation under constraints and feedback; freely
available online.
\end{readinglist}

\chapter{A Geometry of Ideas}

\begin{historicalpreamble}
Modern computing increasingly relies on systems that do not follow
explicit instructions step by step. Instead, they are given a set of
constraints or objectives, and they determine configurations that satisfy
those conditions. Constraint solvers, optimisation algorithms, and
satisfiability systems operate by eliminating impossible configurations
rather than constructing solutions directly.

George Dantzig's simplex algorithm\index[topics]{simplex algorithm}, developed in 1947, showed that the
solution to a linear optimisation problem lies at a vertex of a
geometric region---a corner of the admissible set. The geometry of the
solution space determines where solutions are found. Stephen Cook's 1971
proof that satisfiability is NP-complete\index[topics]{NP-completeness} revealed that the shape of the
constraint space determines the difficulty of search. These developments
confirmed that computation can be understood as geometry: finding stable
points in a landscape defined by constraints.
\end{historicalpreamble}

\begin{historicalfigures}
George Dantzig\index[topics]{Dantzig, George} (1914--2005) developed linear programming and the simplex
method. Stephen Cook\index[topics]{Cook, Stephen} (1939--) proved the NP-completeness of Boolean
satisfiability, founding computational complexity theory. Modern SAT
solver research, building on the work of many contributors since the
1990s, has made constraint-based computation practical at industrial
scale.
\end{historicalfigures}

\begin{paradigmtimeline}
1947: Dantzig's simplex algorithm. \quad
1971: Cook's NP-completeness theorem. \quad
1990s: modern SAT solvers developed. \quad
2000s: constraint-based AI planning and scheduling systems.
\end{paradigmtimeline}

% ── Chapter Figure ────────────────────────────────────────────────────
\begin{figure}[h]
\centering
\begin{tikzpicture}[scale=0.97]
  % Funnel left wall
  \draw[cfig] (-3.2,2.6) .. controls (-2,1.8) and (-1,1.0) .. (0,0.2);
  % Funnel right wall
  \draw[cfig] ( 3.2,2.6) .. controls ( 2,1.8) and ( 1,1.0) .. (0,0.2);
  % Inner funnel lines
  \draw[cfig,gray!50] (-2.2,2.9) .. controls (-1.4,2.1) and (-0.6,1.3) .. (0,0.2);
  \draw[cfig,gray!50] ( 2.2,2.9) .. controls ( 1.4,2.1) and ( 0.6,1.3) .. (0,0.2);
  % Starting points and flow arrows
  \foreach \x/\y in {-2.9/2.65,-2.1/2.85,-0.95/2.35,0/3.05,1.0/2.45,2.1/2.85,2.9/2.65} {
    \fill (\x,\y) circle (1.5pt);
    \draw[->,thin,gray!65]
      (\x,\y) .. controls (\x*0.55,1.65) and (\x*0.18,0.82) .. (0,0.27);
  }
  % Attractor
  \fill[orange!80] (0,0.2) circle (2.5pt);
  \node[orange!85!black] at (0.65,-0.25) {stable attractor};
  \node at (0,3.45) {many initial configurations};
\end{tikzpicture}
\caption{When the constraint landscape has a single stable minimum, every trajectory
in its basin is drawn toward the same attractor. The solution feels inevitable because
every perturbation away from it increases constraint violation.}
\end{figure}

\section{The Shape of Solution Spaces}

We have described constraints as filters on configuration spaces. We can
go further and think about the \emph{geometry} of the admissible set
$\A(\C)$.

\begin{definition}[Constraint Manifold]\index[vocab]{constraint manifold}
When $\A(\C)$ is a smooth subset of a higher-dimensional space $X$, it
is called the \textbf{constraint manifold}. Its curvature measures how
sensitively admissibility responds to perturbations.
\end{definition}

Regions of low curvature are stable: small perturbations remain within
$\A(\C)$. Regions of high curvature are unstable: small perturbations may
leave $\A(\C)$ entirely.

\section{Attractors and Stability}

\begin{definition}[Constraint-Stable Region]\index[vocab]{constraint-stable region}
A region $U \subset \A(\C)$ is \textbf{constraint-stable} if perturbations
originating within $U$ remain within $\A(\C)$ without requiring large
corrective forces.
\end{definition}

The trajectory of a constraint-preserving system is drawn toward stable
regions. This is why certain patterns recur: they are the stable fixed
points of the dynamics, the configurations that resist perturbation.

\section{When Ideas Feel Inevitable}

The feeling of inevitability that accompanies a great idea is the
experiential signal of arriving at a low-curvature region of the constraint
manifold. The idea feels necessary because it is: any small perturbation
from it violates one or more constraints, and the constraint manifold
curves away from those violations.

\begin{principle}[Geometric Characterization of Insight]\index[topics]{insight, geometric characterisation}
An insight corresponds to the recognition of a low-curvature region of
the constraint manifold. The feeling of inevitability is the absence of
nearby violations: every direction from the solution leads quickly into
inadmissibility.
\end{principle}

% ── Formal Correspondence ─────────────────────────────────────────────
\section*{Formal Correspondence}
\begin{formalcorrespondence}
A system seeks configurations minimising constraint violation. Let
$E(x)$ measure constraint violation. Solutions correspond to minima:
\[
  x^* = \arg\min_{x \in \Gamma} E(x).
\]
Closure occurs when $E(x^*) = 0$: all constraints are satisfied
simultaneously and the solution is stable. Low-curvature regions of
$\A(\C)$ correspond to attractors: the constraint manifold curves away
from all nearby inadmissible directions, making the solution feel
necessary. This is the analytic restatement of the geometric principle
above.
\end{formalcorrespondence}

% ── Computational Aside ───────────────────────────────────────────────
\section*{Computational Perspective}
\begin{computationalaside}
Constraint solvers do not search for answers in a constructive sense.
They eliminate possibilities until only admissible configurations remain.
A solution is not selected but forced by the constraints. Strong ideas
exhibit this property: they are not chosen among many but arise as the
only configuration compatible with all conditions simultaneously.
\end{computationalaside}

% ── Worked Example ────────────────────────────────────────────────────
\section*{Worked Example}
Consider the function $E(x) = (x-3)^2(x-7)^2$. Minima occur at $x = 3$
and $x = 7$: these are stable configurations, attractors in the
landscape. If an additional constraint removes one minimum, the system
collapses to the remaining attractor. An idea feels inevitable when the
constraint landscape admits only one stable configuration. The feeling
of necessity is geometric: every direction away from the solution
increases $E$.

\begin{exercise}
List several conditions that must all be satisfied at once---for
example, planning a schedule. Instead of choosing directly, eliminate
options that violate any condition. Observe how the solution emerges
from elimination rather than selection.
\end{exercise}

% ── Spherepop Comparison ──────────────────────────────────────────────
\section*{Spherepop Comparison}
\begin{spherepopcomparison}
A geometry of ideas in Spherepop is a geometry of nested reductions.
Some regions are unstable and collapse away; others behave as attractors.

\textbf{BNF fragment.}
\begin{verbatim}
<landscape> ::= <region-set>
<move>      ::= <region> "=>" <region>
<attractor> ::= "fix(" <region> ")"
\end{verbatim}

An idea feels inevitable when the Spherepop reduction landscape leaves
only one stable attractor after all incompatible nested regions have
been collapsed. This is the grammar-level expression of the geometric
characterisation of insight stated in the principle above.
\end{spherepopcomparison}

\begin{umldiagram}
  \node[box] (land)  {Constraint\\Landscape};
  \node[box, right=2.4cm of land]  (flow) {Reduction\\Flow};
  \node[box, right=2.4cm of flow]  (attr) {Stable\\Attractor};
  \draw[arrow] (land) -- (flow);
  \draw[arrow] (flow) -- (attr);
\end{umldiagram}

% ── Teacher Notes ─────────────────────────────────────────────────────
\section*{Teacher Notes}

\begin{teachernote}
Students should begin to feel the inevitability of solutions, not just
identify them. This is the chapter where qualitative intuition about
constraint systems starts to acquire a geometric character. Encourage
students to think of the admissible set as a landscape with hills and
valleys, and solutions as the lowest points.

The most common difficulty is students viewing answers as arbitrary
selections from a set of equally valid options. Counter this by showing
systems with multiple potential solutions, then adding constraints one
at a time until only one remains. The student should experience the
landscape narrowing around a single stable point.

Watch for students who do not recognise stability. Ask: if I perturb
this solution slightly---change one condition by a small amount---does
it remain admissible? If yes, it is stable. If small perturbations
immediately violate a constraint, the solution is sitting in a sharp
valley---highly stable but fragile. Both are geometrically meaningful
and worth discussing.
\end{teachernote}

\begin{universitybridge}
Energy landscapes and potential surfaces. Dynamical systems attractors
and Lyapunov stability. Optimisation surfaces and saddle points.
Phase space and Poincar\'e's qualitative dynamics.
\end{universitybridge}

\section*{Further Reading}

\begin{readinglist}
Steven Strogatz, \textit{Nonlinear Dynamics and Chaos} --- the most
accessible university-level introduction to attractors and stability;
recommended for motivated students.

James Gleick, \textit{Chaos} --- accessible narrative introduction to
dynamical systems and attractors.

Vladimir Arnold, \textit{Ordinary Differential Equations} --- formal
university treatment of dynamical systems.
\end{readinglist}
% ════════════════════════════════════════════════════════════════════════

\chapter{Deliberate Constraint Accumulation}

\begin{historicalpreamble}
The history of computing is also a history of deliberate constraint
accumulation. Each paradigm shift---from machine code to assembly, from
assembly to structured programming, from procedural to object-oriented,
from sequential to distributed---was not simply the addition of new
capabilities but the imposition of new constraints that made previously
possible errors impossible. Structured programming forbids arbitrary
jumps. Type systems forbid category errors. Memory-safe languages forbid
certain classes of pointer arithmetic. Each restriction narrows the
admissible space of programs, and in doing so makes the programs that
remain more reliably correct.

This is the formal content of software engineering as a discipline: not
the accumulation of tools, but the deliberate accumulation of constraints
that force programs toward admissible configurations. The same logic
applies to thought. Deliberate constraint accumulation is what separates
disciplined inquiry from undisciplined speculation.
\end{historicalpreamble}

\begin{historicalfigures}
Edsger Dijkstra argued that programming is a discipline of constraint:
programs should be constructed so that correctness is provable by
design. Tony Hoare developed the formal precondition--postcondition
framework that makes constraint accumulation explicit. Barbara Liskov's
substitution principle formalised the constraints that object hierarchies
must satisfy to remain coherent.
\end{historicalfigures}

\begin{paradigmtimeline}
1968: Dijkstra's letter on ``goto considered harmful.'' \quad
1969: Hoare logic published. \quad
1980s: type theory and typed functional languages. \quad
2000s: dependent types and proof assistants (Coq, Lean).
\end{paradigmtimeline}

% ── Chapter Figure ────────────────────────────────────────────────────
\begin{figure}[h]
\centering
\begin{tikzpicture}[scale=0.97]
  % Layer 1 (widest)
  \fill[blue!14] (0,0) -- (6,0) -- (5.2,1.0) -- (0.8,1.0) -- cycle;
  \draw[cfig,blue!55!black] (0,0) -- (6,0) -- (5.2,1.0) -- (0.8,1.0) -- cycle;
  % Layer 2
  \fill[green!14] (1.0,1.2) -- (5.0,1.2) -- (4.4,2.0) -- (1.6,2.0) -- cycle;
  \draw[cfig,green!50!black] (1.0,1.2) -- (5.0,1.2) -- (4.4,2.0) -- (1.6,2.0) -- cycle;
  % Layer 3
  \fill[orange!18] (2.0,2.2) -- (4.0,2.2) -- (3.6,2.9) -- (2.4,2.9) -- cycle;
  \draw[cfig,orange!65!black] (2.0,2.2) -- (4.0,2.2) -- (3.6,2.9) -- (2.4,2.9) -- cycle;
  % Final stable point
  \fill[red!70] (3.0,3.15) circle (2.4pt);
  % Layer labels
  \node[blue!65!black]   at (6.65,0.5)  {$C_1$};
  \node[green!55!black]  at (5.65,1.65) {$C_2$};
  \node[orange!70!black] at (4.55,2.6)  {$C_3$};
  \node[red!70!black]    at (3.75,3.3)  {solution};
  % Brace annotation
  \draw[<->,thin,gray!55] (-0.4,0) -- (-0.4,3.15);
  \node[gray!65!black,rotate=90] at (-0.85,1.55) {accumulation};
\end{tikzpicture}
\caption{Each constraint layer narrows the admissible space. Deliberately added constraints
converge toward a single stable result that could not have been reached by any one layer alone.}
\end{figure}

\section{A Practical Framework}

The formal picture suggests a practical strategy for accelerating the
formation of ideas.

\begin{enumerate}
  \item \textbf{Increase constraint density.} Learn from multiple domains,
        not just one. Each domain imposes constraints that reduce the
        admissible space for all domains simultaneously.
  \item \textbf{Seek overlap.} Find places where constraints from different
        domains apply to the same configuration space. These are the regions
        where accumulation accelerates.
  \item \textbf{Test for global consistency.} When a candidate idea appears,
        do not only test it within its domain of origin. Apply constraints
        from every other relevant domain. If it survives, it is a candidate
        for a real invariant.
  \item \textbf{Diagnose failure modes.} When an idea breaks down, identify
        whether the failure is local (fragmentation: one piece does not fit)
        or global (obstruction: the pieces cannot be reconciled). The
        diagnosis determines the repair.
\end{enumerate}

\section{The Role of Formalism}

Mathematical formalism is a constraint-imposing technology. By requiring
that ideas be expressed in precise, compositional language, formalism
eliminates imprecision that allows contradictions to hide.

This is why formalism is useful even in fields that are not primarily
mathematical. The discipline of writing a precise definition, or of
proving a claim rather than asserting it, forces confrontation with
the constraints that informal language evades.

\section{Open Problems}

\begin{openproblem}
How can the formation of scientific theories be modeled as constraint
accumulation across the domains of experimental evidence, mathematical
consistency, and theoretical coherence? What does the constraint manifold
of a scientific paradigm look like, and how does a paradigm shift correspond
to a change in its topology?
\end{openproblem}

\begin{openproblem}
What is the relationship between the structural entropy defined here and
Kolmogorov complexity? Both measure a form of compressibility. Are they
equivalent under any natural mapping?
\end{openproblem}

\begin{openproblem}
Can the conditions for structural emergence (Chapter~2) be given a
quantitative threshold? That is, is there a formal criterion for when
a system is ``constraint-dense enough'' to reliably generate insight?
\end{openproblem}

% ── Formal Correspondence ─────────────────────────────────────────────
\section*{Formal Correspondence}
\begin{formalcorrespondence}
Deliberate constraint accumulation corresponds to constructing the
admissible set $\Gamma_{\mathcal{C}}$ incrementally and intentionally.
At each step, add a constraint $C_{n+1}$ and recompute:
\[
  \Gamma_{\mathcal{C} \cup \{C_{n+1}\}} = \Gamma_{\mathcal{C}} \cap \Gamma_{n+1}.
\]
The process terminates in closure when $|\Gamma_{\mathcal{C}}|$ is
small enough to identify a solution. Formalism is a constraint-imposing
technology: precise language forces ideas to be stated in ways that make
their constraint satisfaction testable. This expresses the practical
framework above in the language of the admissible set.
\end{formalcorrespondence}

% ── Computational Aside ───────────────────────────────────────────────
\section*{Computational Perspective}
\begin{computationalaside}
Each programming paradigm shift in history was the imposition of new
constraints that made previously possible errors impossible. Structured
programming forbids arbitrary jumps; type systems forbid category errors;
memory-safe languages forbid certain pointer operations. Each restriction
narrows the admissible space and, in doing so, makes the programs that
remain more reliably correct. This is deliberate constraint accumulation
applied to computation itself.
\end{computationalaside}

% ── Worked Example ────────────────────────────────────────────────────
\section*{Worked Example}
Suppose we seek a number satisfying: $x > 10$; $x$ is even; $x$ is
divisible by 3. Accumulated constraints yield $x \in \{12, 18, 24,\dots\}$.
Adding a refinement $x < 20$ gives $x \in \{12, 18\}$. Deliberate
constraint accumulation narrows the solution space in a controlled way:
each constraint is chosen because it eliminates a class of inadmissible
possibilities, not because it directly names the answer.

\begin{exercise}
Choose a decision you need to make. List five constraints it must
satisfy. Intersect them one at a time. Observe how the admissible space
shrinks with each addition, and whether the constraints eventually force
a unique option.
\end{exercise}

% ── Spherepop Comparison ──────────────────────────────────────────────
\section*{Spherepop Comparison}
\begin{spherepopcomparison}
Deliberate constraint accumulation in Spherepop means intentionally
constructing nested regions so that invalid configurations collapse
early and only strong outer forms remain.

\textbf{BNF fragment.}
\begin{verbatim}
<build>     ::= "add(" <constraint> "," <region> ")"
<refine>    ::= "refine(" <region> ")"
<stabilize> ::= "stable(" <region> ")"
\end{verbatim}

Building Spherepop structures that increase density, preserve closure,
and make projection possible is the grammar-level expression of the
practical framework described in this chapter.
\end{spherepopcomparison}

\begin{umldiagram}
  \node[box] (raw)    {Initial\\Region Space};
  \node[box, right=2.4cm of raw]    (ref)  {Deliberate\\Refinement};
  \node[box, right=2.4cm of ref]    (stab) {Stable\\Configuration};
  \draw[arrow] (raw) -- (ref);
  \draw[arrow] (ref) -- (stab);
\end{umldiagram}

% ── Teacher Notes ─────────────────────────────────────────────────────
\section*{Teacher Notes}

\begin{teachernote}
Students should learn to construct problems intentionally, not just solve
them. The shift from passive problem-solving to active problem design is
the central pedagogical aim of this chapter. Ask students to create a
constraint system that has exactly one solution, then ask them to justify
why no other configuration is admissible. This reversal---designing
constraints to force an outcome---is the most powerful exercise in the
book.

The most common difficulty is a passive problem-solving mindset. Students
who have spent years solving pre-formed problems often find it
disorienting to design the constraints themselves. Start with small,
fully specified examples: ``create three constraints on numbers from 1 to
20 such that exactly one number satisfies all three.''

Watch for overly weak constraints (many solutions remain), contradictory
constraints (no solutions remain), and redundant constraints (later
constraints add no new elimination). All three are instructive failure
modes. Ask students to diagnose which type of failure they have produced.
\end{teachernote}

\begin{universitybridge}
Formal methods in software engineering. Constraint programming
languages (Prolog, MiniZinc). Axiomatic system design. Engineering
design principles and requirements specification.
\end{universitybridge}

\section*{Further Reading}

\begin{readinglist}
Donald Knuth, \textit{The Art of Computer Programming} --- systematic
construction of algorithms through deliberate structural design; a
life's work on constraint accumulation applied to computation.

Niklaus Wirth, \textit{Algorithms + Data Structures = Programs} ---
structured problem design; concise and influential.

Edsger Dijkstra, \textit{A Discipline of Programming} --- argues for
constructing programs so that correctness is provable by design; the
best model of deliberate constraint accumulation in software.
\end{readinglist}

\chapter{Conclusion: Structure Is Primary}

\begin{historicalpreamble}
The history of computation is, at its deepest level, the progressive
discovery that structure precedes content. Early machines computed
numbers; later systems manipulated symbols; later still they processed
relationships. At each stage, what the machine operated on became more
abstract---and what mattered more was the structure of the operations
than the nature of the objects they operated on.

Alan Turing showed that a machine capable of simulating any other
machine need not know what those machines compute---only how they are
structured. Alonzo Church showed that functions, not numbers, are the
primitive objects of computation. John von Neumann showed that programs
and data share the same representational space. Each of these insights
is a version of the same claim: structure is primary. The objects are
secondary. What persists across different representations, different
domains, different computational models, is invariant structure---the
same thing this book has been about from the first page.
\end{historicalpreamble}

\begin{historicalfigures}
Alan Turing\index[topics]{Turing, Alan}'s 1936 paper on computable numbers established that any
effective procedure can be performed by a machine that reads and writes
symbols according to a finite set of rules---structure, not substance.
Alonzo Church's lambda calculus showed that computation is the
transformation of expressions under rules of substitution. John von
Neumann's stored-program architecture unified program and data, making
explicit that the distinction between structure and content is not fixed
but depends on the level of description.
\end{historicalfigures}

\begin{paradigmtimeline}
1936: Turing's paper on computable numbers; Church's lambda calculus. \quad
1945: von Neumann architecture draft. \quad
1960s--70s: programming language theory unifies the models. \quad
1990s--present: type theory, proof assistants, and category theory
formalise the unity of structure across mathematics and computation.
\end{paradigmtimeline}

% ── Chapter Figure ────────────────────────────────────────────────────
\begin{figure}[h]
\centering
\begin{tikzpicture}[scale=0.97]
  % Three surface forms (top row)
  \draw[cfig] (0,3.1) circle (0.85);
  \node at (0,4.2) {form 1};
  \draw[cfig] (3.2,2.3) -- (4.45,3.95) -- (5.2,2.45) -- cycle;
  \node at (4.1,4.2) {form 2};
  \draw[cfig] (7.1,2.3) rectangle (8.75,3.85);
  \node at (7.9,4.2) {form 3};
  % Projection arrows
  \draw[->,thick,gray!65] (0.0,2.2)  to[bend left=10]  (2.8,0.75);
  \draw[->,thick,gray!65] (4.1,2.15) -- (4.1,0.75);
  \draw[->,thick,gray!65] (7.9,2.2)  to[bend right=10] (5.4,0.75);
  % Invariant skeleton (bottom)
  \draw[very thick] (2.8,0) -- (5.4,0);
  \draw[very thick] (4.1,-0.85) -- (4.1,0.85);
  \fill (4.1,0) circle (2.2pt);
  \node at (4.1,-1.3) {shared invariant structure};
  % Labels on arrows
  \node[gray!65!black] at (1.5,1.7) {$\pi_1$};
  \node[gray!65!black] at (4.5,1.4) {$\pi_2$};
  \node[gray!65!black] at (6.7,1.7) {$\pi_3$};
\end{tikzpicture}
\caption{Three surface representations---circle, triangle, rectangle---each project onto
the same underlying invariant structure. The representation changes with the domain;
the structure does not.}
\end{figure}

\section{The Core Claim}

This book has developed and applied a single claim:

\begin{principle}[Primacy of Structure]\index[topics]{primacy of structure}
Ideas are not created; they are recognized. Recognition occurs when a
cognitive system has accumulated enough constraints from enough domains
to reduce the admissible configuration space to a unique element. The
element was always the only admissible one; accumulation is what makes
it visible.
\end{principle}

\section{Consequences}

If this is right, several things follow.

Originality is not the same as novelty. The most original ideas are those
that satisfy the most demanding set of constraints---constraints that few
people have accumulated simultaneously. Novelty for its own sake imposes
few constraints and is therefore easy to produce but structurally shallow.

Expertise is not knowledge retention. It is constraint accumulation and
successful relegation. An expert has internalized many constraints and has
compressed the admissible structure of their domain into rapidly deployable
form.

Failure is not waste. Every constraint failure---every time a candidate
idea runs into a constraint that eliminates it---is information. It reduces
the admissible set and moves the system closer to the true configuration.

Formalism is not decoration. Formal languages are constraint-imposing
technologies. They force ideas to be precise enough to be tested.

\section{Where to Go From Here}

The framework sketched here---constraint accumulation, structural entropy,
projection, admissibility, and functorial correspondence---has formal
foundations in several mathematical disciplines: set theory, topology,
information theory, and category theory. Each of these offers tools that
sharpen and extend what is developed here.

The most direct extensions are:
\begin{itemize}
  \item \textbf{Category theory:} provides the formal language for
        functorial correspondence and structure-preserving maps.
  \item \textbf{Sheaf theory:} formalizes the local-global distinction and
        the conditions under which local sections extend to global ones.
  \item \textbf{Information theory:} grounds the entropy formalism in a
        rigorous probabilistic framework.
  \item \textbf{Differential geometry:} provides the tools to study the
        constraint manifold, its curvature, and its attractors.
\end{itemize}

None of these require more prerequisite knowledge than a strong high school
mathematics background to begin. They are the natural next layer of
constraint accumulation for someone who has read this book.

% ── Formal Correspondence ─────────────────────────────────────────────
\section*{Formal Correspondence}
\begin{formalcorrespondence}
Two representations are structurally equivalent when they reduce to the
same invariant under the constraint set:
\[
  \mathrm{form}_1 \sim \mathrm{form}_2 \iff \pi(\mathrm{form}_1)
  = \pi(\mathrm{form}_2).
\]
Structure is primary because what survives across distinct representations
is not the surface form but the invariant pattern of admissible reduction.
Different descriptions of the same idea are different coordinate systems
on the same constraint manifold. The admissible set $\A(\C)$ is the
invariant; the representations are projections of it. This is the formal
content of the principle of primacy of structure.
\end{formalcorrespondence}

% ── Computational Aside ───────────────────────────────────────────────
\section*{Computational Perspective}
\begin{computationalaside}
The history of computation is the progressive discovery that structure
precedes content. Turing showed that a universal machine need not know
what other machines compute---only how they are structured. Church showed
that functions, not numbers, are the primitive objects of computation.
Von Neumann showed that programs and data share the same representational
space. Each insight is a version of the same claim: structure is primary,
objects are secondary, and what persists across representations is
invariant structure.
\end{computationalaside}

% ── Worked Example ────────────────────────────────────────────────────
\section*{Worked Example}
Consider two different descriptions: $(2 + 3) \times 4$ and $5 \times 4$.
Though syntactically different, both reduce to 20. The underlying
structure determines the result, not the surface form. Different
representations collapse to the same invariant outcome. The entire
framework of this book is a single extended instance of this observation:
constraints, entropy, abstraction, functors, and attractors are all
different coordinate systems on the same invariant structure.

% ── Spherepop Comparison ──────────────────────────────────────────────
\section*{Spherepop Comparison}
\begin{spherepopcomparison}
Spherepop ends where this book ends: different surface descriptions may
correspond to the same stabilised structure once irrelevant inner
distinctions have been reduced.

\textbf{BNF fragment.}
\begin{verbatim}
<form1>     ::= <region>
<form2>     ::= <region>
<invariant> ::= "same(" <form1> "," <form2> ")"
\end{verbatim}

What survives across distinct Spherepop representations is not the
surface grammar but the invariant pattern of admissible reduction. This
is the grammar-level expression of the primacy of structure---the same
claim that opens and closes this book.
\end{spherepopcomparison}

\begin{umldiagram}
  \node[box] (f1)  {Representation\\One};
  \node[box, right=3.2cm of f1]  (f2)  {Representation\\Two};
  \node[box, below=1.8cm of $(f1)!0.5!(f2)$] (inv) {Shared\\Invariant Structure};
  \draw[arrow] (f1) -- (inv);
  \draw[arrow] (f2) -- (inv);
\end{umldiagram}

% ── Teacher Notes ─────────────────────────────────────────────────────
\section*{Teacher Notes}

\begin{teachernote}
Students should arrive at this chapter able to recognise invariance
across representations without being told what to look for. The most
common difficulty is confusing surface form with meaning: students who
have spent years learning to read specific notations often mistake the
notation for the thing it describes. Design exercises that present the
same structure in three radically different forms and ask students to
identify what is preserved across all three.

This chapter is also a moment for reflection. Ask students to look back
at Chapter~1 and describe what they understand now that they did not
understand then. The answer should not be a list of new topics, but a
description of the same idea seen with greater clarity. If a student
can articulate that, the book has done its work.

Watch for students who focus on surface differences. The question to
ask is always: if I change the representation, what changes in the
structure, and what does not? The invariants are what does not change.
Those are what the book has been about.
\end{teachernote}

\begin{universitybridge}
Isomorphism in abstract algebra. Invariant theory. Abstract algebra
and group theory. Theoretical computer science and formal language
equivalences.
\end{universitybridge}

\section*{Further Reading}

\begin{readinglist}
Hermann Weyl, \textit{Symmetry} --- explores invariance across
mathematics, art, and science; accessible and beautifully written.

Eugene Wigner, \textit{The Unreasonable Effectiveness of Mathematics
in the Natural Sciences} (1960) --- a short, famous essay on why
mathematical structure appears across domains; worth reading in full.

Terence Tao, \textit{An Introduction to Measure Theory} (selected
ideas) --- an advanced example of structure taking precedence over
representation; the preface is accessible and illuminating.
\end{readinglist}

% ════════════════════════════════════════════════════════════════════════
\backmatter
% ════════════════════════════════════════════════════════════════════════

% ════════════════════════════════════════════════════════════════════════
\chapter*{What You Are Now Prepared to Study}
\addcontentsline{toc}{chapter}{What You Are Now Prepared to Study}
% ════════════════════════════════════════════════════════════════════════

\begin{quote}
Throughout this text, you have encountered the same structure in many
different forms.

You have seen problems where answers did not need to be invented, but
emerged as the only configurations that satisfied all conditions. You
have seen that simplifying a system does not remove structure, but
reveals what remains unchanged. You have seen that systems can appear
correct in parts while failing as a whole, and that true understanding
requires consistency across all perspectives.

At each stage, these ideas were presented through examples, history,
computation, and formal reasoning. None of these were separate topics.
They were different ways of describing the same underlying process.

You are now prepared to recognise that process in other fields.
\end{quote}

\section*{Mathematics}

\begin{quote}
In mathematics, you will encounter systems defined by constraints and
relationships. Solutions are not arbitrary---they are determined by the
structure of those relationships.

You will study objects such as sets, functions, spaces, and algebraic
systems. What matters is not the objects themselves, but the invariants
they preserve and the transformations that leave them unchanged.

You will also encounter situations where local consistency does not
guarantee global coherence, and where different representations describe
the same underlying structure.

The patterns you have already seen---constraint accumulation, reduction,
and stabilisation---will appear in more formal and precise forms.
\end{quote}

\section*{Computer Science}

\begin{quote}
In computer science, you will study computation as a process of
transformation.

Some systems will operate step by step, applying rules in sequence.
Others will reduce expressions until they reach stable forms. Still
others will search for configurations that satisfy all constraints at
once.

You will encounter ideas such as algorithms, data structures,
programming languages, and distributed systems. In each case, the
central question is not how to produce output, but how to maintain
consistency across all conditions imposed on the system.

What you have learned here---that computation is the structured
resolution of constraints---will apply across all of these areas.
\end{quote}

\section*{Physics and Information}

\begin{quote}
In physics, systems are described by laws that restrict what can occur.
The behaviour of a system is determined by the configurations that
satisfy those laws.

You will encounter concepts such as energy, entropy, and conservation.
Entropy measures how many configurations are possible. Energy often
measures how constrained a system is.

You will also see that processes are not always reversible, and that
history can be compressed into states that cannot be fully recovered.

The same principles you have studied---multiplicity, compression, and
irreversibility---appear here as fundamental laws of nature.
\end{quote}

\section*{Logic and Formal Systems}

\begin{quote}
In logic, you will study systems of statements and rules.

A system is consistent if all constraints can be satisfied together.
It is inconsistent if no configuration satisfies them all.

You will learn to distinguish between local validity and global
consistency, and to recognise when a system cannot be completed without
contradiction.

These are not new ideas. They are the formal expression of patterns you
have already encountered.
\end{quote}

\section*{Cross-Domain Thinking}

\begin{quote}
As you continue, you will notice that the same structures appear in
different fields.

A system in mathematics may resemble a system in physics. A process in
computation may resemble a process in reasoning. A structure in one
domain may be translated into another while preserving its essential
relationships.

Recognising these correspondences is not a matter of analogy, but of
structure. What is preserved across domains is not the surface form,
but the relationships that define the system.

You are now prepared to recognise when two different problems are, in
fact, the same problem expressed in different terms.
\end{quote}

\section*{What Changes Now}

\begin{quote}
At the beginning of this text, problems may have appeared as things to
be solved.

At this point, problems should begin to appear as systems to be
understood.

You may find that answers feel less like choices and more like
consequences. You may notice that certain solutions seem inevitable
once all conditions are taken into account.

This is not because the problems have become easier, but because their
structure has become visible.

The purpose of this text has not been to teach you a collection of
methods. It has been to make that structure visible.

From here, you will encounter more precise languages, more complex
systems, and more demanding problems. But the underlying patterns will
remain the same.

You are not beginning something new. You are continuing the same
process, with greater clarity.
\end{quote}

% ════════════════════════════════════════════════════════════════════════
\chapter*{Pathways Beyond This Book}
\addcontentsline{toc}{chapter}{Pathways Beyond This Book}
% ════════════════════════════════════════════════════════════════════════

The reading lists at the end of each chapter point outward from
individual topics. This section groups those recommendations into
three coherent trajectories. A student who follows any one of these
paths will find that the concepts from this book recur---in more
precise forms, with more powerful tools, at greater depth.

\section*{The Mathematics Pathway}

This path moves from the informal set theory and constraint language of
this book into algebra, topology, and category theory. It is the
trajectory toward pure mathematics.

\begin{readinglist}
\textit{Starting point:} G.~P\'olya, \textit{How to Solve It}. Builds
the habit of constraint-based reasoning informally.

\textit{First year:} F.~W.~Lawvere and S.~Schanuel,
\textit{Conceptual Mathematics}. Introduces category theory with no
prerequisites; formalises functors, structure-preserving maps, and
invariants.

\textit{Second year:} M.~Artin, \textit{Algebra}. Standard
first-course abstract algebra; groups, rings, and fields as constraint
systems on sets of operations.

\textit{Third year:} A.~Hatcher, \textit{Algebraic Topology}.
Formalises local-global consistency, obstruction, and cohomology.

\textit{Ongoing:} H.~Weyl, \textit{Symmetry}; E.~Wigner,
\textit{The Unreasonable Effectiveness of Mathematics}. Both repay
rereading at every level.
\end{readinglist}

\section*{The Computer Science Pathway}

This path moves from the computational paradigm history in each chapter
into algorithm theory, programming languages, and formal verification.

\begin{readinglist}
\textit{Starting point:} H.~Abelson and G.~J.~Sussman,
\textit{Structure and Interpretation of Computer Programs} (freely
available online). The best single book for seeing abstraction,
reduction, and constraint as computational primitives.

\textit{First year:} M.~Sipser, \textit{Introduction to the Theory of
Computation}. Covers automata, formal languages, computability, and
complexity; places computation in a structural framework.

\textit{Second year:} T.~Cormen et al., \textit{Introduction to
Algorithms}. Constraint-based algorithm design at scale.

\textit{Third year:} B.~Pierce, \textit{Types and Programming
Languages}. Type theory as a system of constraints on programs;
connects directly to quotient structures and abstraction.

\textit{Ongoing:} E.~Dijkstra, \textit{A Discipline of Programming};
D.~Knuth, \textit{The Art of Computer Programming}.
\end{readinglist}

\section*{The Physics and Information Pathway}

This path moves from the entropy and compression chapters into
thermodynamics, information theory, and the physics of computation.

\begin{readinglist}
\textit{Starting point:} C.~Shannon, \textit{A Mathematical Theory of
Communication} (1948). Short, precise, and foundational; the first few
sections are accessible after this book.

\textit{First year:} D.~MacKay, \textit{Information Theory, Inference,
and Learning Algorithms} (freely available online). Covers entropy,
compression, and inference in a unified framework.

\textit{Second year:} R.~Feynman, \textit{The Feynman Lectures on
Physics} (Vol.~I, thermodynamics sections). Connects entropy and
constraint to physical systems.

\textit{Third year:} T.~Cover and J.~Thomas, \textit{Elements of
Information Theory}. The standard university treatment.

\textit{Ongoing:} C.~H.~Bennett, \textit{Logical Reversibility of
Computation}; S.~Strogatz, \textit{Nonlinear Dynamics and Chaos}.
\end{readinglist}

% ════════════════════════════════════════════════════════════════════════
\chapter*{Glossary}
\addcontentsline{toc}{chapter}{Glossary}
% ════════════════════════════════════════════════════════════════════════

\textbf{Abstraction.} A map from a space of configurations to a simpler
space that preserves admissibility while identifying configurations that
are equivalent under the constraint set. Abstraction hides completed
internal computation, not ongoing computation.

\textbf{Admissible set.} The set of all configurations that satisfy
every constraint in a given constraint set. Written $\mathcal{A}(\mathcal{C})$.
As constraints accumulate, the admissible set shrinks.

\textbf{Attractor.} A stable configuration in a dynamical system or
constraint landscape. Small perturbations from an attractor return to it.
An attractor corresponds to an idea that feels inevitable because every
nearby alternative violates some constraint.

\textbf{Closure.} The condition that the admissible set is non-empty:
$\mathcal{A}(\mathcal{C}) \neq \emptyset$. Closure is the precondition for any
solution to exist.

\textbf{Compression.} A map that reduces the size of a representation
while preserving relevant structure. Constraint-preserving compression
identifies configurations only when they are genuinely equivalent under
the constraint set.

\textbf{Constraint.} A condition on a configuration space that
eliminates configurations that do not satisfy it. A constraint is a
filter, not a recipe.

\textbf{Constraint density.} The degree to which constraints in a
system interact non-trivially, producing an admissible set smaller than
the intersection of each constraint taken alone.

\textbf{Entropy.} The logarithm of multiplicity: $S = \log \mu$.
Entropy measures how many configurations are consistent with the
current state of knowledge. It is not disorder; it is underdetermination.

\textbf{Fixed point.} A configuration $x$ such that $T(x) = x$ under
a transformation $T$. Fixed points correspond to stable solutions and
to the normal forms of computational reduction.

\textbf{Formal correspondence.} A precise restatement of a conceptual
claim in mathematical language. In this book, formal correspondences
are presented as clarifications of ideas already understood, not as
new claims.

\textbf{Functor.} A structure-preserving map between categories
(domains with compositional rules). A functor translates the grammar
of relationships from one domain into another without distorting it.

\textbf{Global consistency.} The condition that a configuration
satisfies all constraints across the entire domain, not just locally.
Global consistency does not follow from local consistency.

\textbf{Hallucination.} A configuration that appears globally
consistent under a restricted projection but does not satisfy the full
constraint set. Produced by overcompressive abstraction.

\textbf{Invariant.} A property preserved by a given transformation or
abstraction. Invariants are what remain the same when representations
change.

\textbf{Local consistency.} The condition that a configuration satisfies
all constraints within a restricted region. Necessary but not sufficient
for global consistency.

\textbf{Multiplicity.} The number of configurations in a full space
that map to the same configuration in an abstracted or compressed
representation. Written $\mu(y)$.

\textbf{Obstruction.} A condition in which two locally consistent
structures cannot be combined into a globally consistent whole. An
obstruction is not a failure of effort but a structural fact.

\textbf{Overcompression.} An abstraction that identifies configurations
that are distinguishable under the constraint set, discarding relevant
distinctions and producing false simplicity.

\textbf{Projection capacity.} The ability of a system to map its
admissible configurations into a simpler representation while
preserving their essential structure.

\textbf{Relegation.} The process by which a cognitive system compresses
a constraint structure into a latent form that can be retrieved and
applied without active reconstruction. Intuition is relegated
constraint satisfaction.

\textbf{Spherepop Calculus.} A formal system in which structure is
built from nested regions, merges, and collapses. Used in this book
as a grammar-level coordinate system for the same structural ideas
expressed in set theory and category theory.

\textbf{Structure-preserving map.} A function between domains that maps
admissible configurations to admissible configurations and respects
compositional rules. Called a functor in category theory.

\textbf{Temporal compression.} A map from a full history of
configurations to the present state. When temporal compression is lossy,
the process is irreversible: the past cannot be recovered from the
present.

% ════════════════════════════════════════════════════════════════════════
\chapter*{Mathematical Appendix}
\addcontentsline{toc}{chapter}{Mathematical Appendix}
% ════════════════════════════════════════════════════════════════════════

This appendix collects the mathematical background used in the main
text. Each section covers one concept informally, then states it
precisely enough to support the definitions and theorems in the
chapters.

\section*{Sets and Subsets}

A \emph{set} is any collection of objects. We write $x \in A$ to mean
that $x$ belongs to $A$, and $x \notin A$ to mean it does not. We write
$A \subseteq B$ when every element of $A$ also belongs to $B$.

The \emph{intersection} of two sets is
$A \cap B = \{x \mid x \in A \text{ and } x \in B\}$.
The intersection of many sets is $\bigcap_{i=1}^n A_i$, the set of
elements belonging to all of them.

The \emph{empty set} $\emptyset$ contains no elements.

\section*{Functions}

A \emph{function} $f : X \to Y$ assigns to each $x \in X$ exactly one
$f(x) \in Y$. The set $X$ is the domain; $Y$ is the codomain.

A function is \emph{surjective} (onto) if every element of $Y$ is
$f(x)$ for some $x \in X$. It is \emph{injective} (one-to-one) if
$f(x_1) = f(x_2)$ implies $x_1 = x_2$. It is \emph{bijective} if
both.

The \emph{pre-image} of $y \in Y$ is $f^{-1}(y) = \{x \in X \mid
f(x) = y\}$.

\section*{Equivalence Relations}

A relation $\sim$ on a set $X$ is an \emph{equivalence relation} if it
is reflexive ($x \sim x$), symmetric ($x \sim y$ implies $y \sim x$),
and transitive ($x \sim y$ and $y \sim z$ imply $x \sim z$).

An equivalence relation partitions $X$ into disjoint \emph{equivalence
classes}. The \emph{quotient set} $X / {\sim}$ is the set of those
classes. A function $\pi : X \to X/{\sim}$ sending each element to its
class is the \emph{quotient map}. This is the formal structure of
abstraction.

\section*{Logarithms}

The \emph{logarithm base $b$} of $n$, written $\log_b n$, is the
unique real number $e$ such that $b^e = n$. Throughout this book we
use $\log = \log_2$, following the convention of information theory.

Key properties: $\log(mn) = \log m + \log n$; $\log(m/n) = \log m -
\log n$; $\log(m^k) = k \log m$.

If a system has $n$ equally likely configurations, the entropy is
$\log n$ bits. Adding a constraint that halves the admissible set
reduces entropy by exactly $1$ bit.

\section*{Intersections as Constraint Accumulation}

If $\mathcal{C} = \{C_1, \dots, C_n\}$ is a set of constraints, each
defining a subset $\Gamma_i \subseteq \Gamma$, then the admissible set
is $\mathcal{A}(\mathcal{C}) = \bigcap_{i=1}^n \Gamma_i$.

Adding a constraint $C_{n+1}$ replaces $\mathcal{A}(\mathcal{C})$ with
$\mathcal{A}(\mathcal{C}) \cap \Gamma_{n+1} \subseteq
\mathcal{A}(\mathcal{C})$. Constraint accumulation is monotone:
the admissible set never grows when constraints are added.

\section*{A Note on Proofs}

The proofs in this book are sketches: they convey the structure of the
argument without supplying every technical detail. A sketch is a
constraint on what a full proof must look like. It eliminates the
space of possible proofs to a small region; filling in the details
locates the unique proof within that region.

This is itself an instance of the book's central idea.

% ════════════════════════════════════════════════════════════════════════
\chapter*{Computational Appendix}
\addcontentsline{toc}{chapter}{Computational Appendix}
% ════════════════════════════════════════════════════════════════════════

This appendix gives brief, self-contained examples of the computational
ideas discussed in the chapter historical preambles. No programming
experience is required to read them. The goal is to make the paradigm
descriptions concrete.

\section*{Imperative Style: Constraint as Sequential Elimination}

An imperative program applies rules in order. Each rule eliminates
possibilities. The following pseudocode narrows a number:

\begin{verbatim}
candidates = {1, 2, ..., 100}
candidates = {x in candidates : x is even}
candidates = {x in candidates : x > 50}
candidates = {x in candidates : x < 70}
-- candidates is now {52, 54, 56, 58, 60, 62, 64, 66, 68}
\end{verbatim}

Each line is a constraint. Each constraint reduces the set.

\section*{Functional Style: Reduction to Normal Form}

A functional program defines transformations and reduces expressions.

\begin{verbatim}
simplify((x + 0) * 1)
  = simplify(x * 1)      -- x + 0 reduces to x
  = simplify(x)          -- x * 1 reduces to x
  = x                    -- normal form reached
\end{verbatim}

Reduction continues until no rule applies. The result is the normal
form: a fixed point of the reduction system.

\section*{Lazy Evaluation: Deferred Computation}

A lazy system carries unevaluated expressions forward and resolves
them only when required.

\begin{verbatim}
stream = [1, 2, 3, 4, ...]          -- infinite list, not yet evaluated
first_over_100 = first(x in stream : x > 100)
-- evaluates stream only until 101 is reached
-- result: 101
\end{verbatim}

The computation is deferred. The full stream is never evaluated.
Insight occurs when sufficient evaluation has accumulated.

\section*{Modular Abstraction: Interface over Implementation}

A module exposes behaviour without exposing internal structure.

\begin{verbatim}
module Stack:
  push(item)    -- add item to top
  pop()         -- remove and return top item
  is_empty()    -- return true if no items remain

-- Internal implementation (array, linked list, etc.) is hidden.
-- The interface is the constraint. The implementation is irrelevant
-- to any code that uses the module.
\end{verbatim}

Abstraction means the internal structure has been resolved to the point
that it no longer affects external interaction.

\section*{Spherepop BNF: The Full Grammar}

The Spherepop Calculus used throughout this book is defined by the
following grammar. A \emph{region} is either an atom or a nested list
of regions. Operations include merging, collapsing, storing, and
abstracting regions.

\begin{verbatim}
<system>      ::= <region>
<region>      ::= <atom>
               |  "(" <region-list> ")"
<region-list> ::= <region>
               |  <region> <region-list>
<op>          ::= "merge("  <region> "," <region> ")"
               |  "collapse(" <region> ")"
               |  "abs("   <region> ")"
               |  "store("  <region> ")"
               |  "use("   <latent>  ")"
<latent>      ::= "store(" <region> ")"
<constraint>  ::= "allow"  <region>
               |  "forbid" <region>
\end{verbatim}

A region is admissible if it is not forbidden by any active constraint.
A collapsed region is admissible only if its interior was admissible
before collapse. An abstracted region exposes only its outer boundary.

% ════════════════════════════════════════════════════════════════════════
\chapter*{Topic Index}
\addcontentsline{toc}{chapter}{Topic Index}
% ════════════════════════════════════════════════════════════════════════

\begin{quote}\small
\textit{Topics, named theorems and principles, computational paradigms,
historical figures, and cross-domain concepts.}
\end{quote}

\printindex[topics]

% ════════════════════════════════════════════════════════════════════════
\chapter*{Vocabulary Index}
\addcontentsline{toc}{chapter}{Vocabulary Index}
% ════════════════════════════════════════════════════════════════════════

\begin{quote}\small
\textit{Formally defined terms, pointing to the definition environment
where each term is introduced. Subentries indicate specialised uses.}
\end{quote}

\printindex[vocab]

\end{document}
